\documentclass[12pt,a4paper]{article}
\usepackage[utf8]{inputenc}
\usepackage[T1]{fontenc}
\usepackage[ngerman,latin]{babel}
\usepackage{amsmath,amssymb,amsthm}
\usepackage{geometry}
\usepackage{hyperref}
\usepackage{enumitem}
\usepackage{mathrsfs}
\usepackage{longtable}
\usepackage{booktabs}
\usepackage{microtype}
\usepackage{etoolbox}
\geometry{margin=2.5cm}

\theoremstyle{plain}
\newtheorem{satz}{Satz}[section]
\newtheorem{theorem}{Theorem}[section]
\newtheorem{lemma}{Lemma}[section]
\newtheorem{hilfssatz}{Hilfssatz}[section]
\newtheorem{korollar}{Korollar}[section]

\theoremstyle{definition}
\newtheorem{definition}{Definition}[section]

\theoremstyle{remark}
\newtheorem{bemerkung}{Bemerkung}[section]
\newtheorem{beispiel}{Beispiel}[section]

\title{Carl Friedrich Gau\ss{} Werke --- Band 1, Teil 2}
\author{Carl Friedrich Gau\ss{}}
\date{}

\begin{document}
\maketitle


\section*{MODULI QUI SUNT NUMERI PRIMI.}

\setcounter{satz}{49}
\begin{satz}[Fermati Theorema.]
\setcounter{satz}{50}
\addtocounter{satz}{-1}
\textbf{50.}

Quum igitur $\frac{p-1}{t}$ sit integer, sequitur evehendo utramque partem congruentiae $a^t \equiv 1$ ad potestatem exponentis $\frac{p-1}{t}$, $a^{p-1} \equiv 1$, sive $a^{p-1} - 1$ semper per $p$ divisibilis est, quando $p$ est primus ipsum $a$ non metiens.
\end{satz}

Theorema hoc, quod tum propter elegantiam tum propter eximiam utilitatem omni attentione dignum, ab inventore theorema \emph{Fermatianum} appellari solet. Vid. Fermatii Opera Mathem. Tolosae 1619 fol. p.~163. Demonstrationem inventor non adiecit, quam tamen in potestate sua esse professus est. Ill. Euler primus demonstrationem publici iuris fecit, in diss. cui titulus \emph{Theorematum quorundam ad numeros primos spectantium demonstratio}, Comm. Acad. Petrop. T.~VIII\footnote{In comment. anteriore vir summus ad scopum nondum pervenerat. Comm. Petr. T.~VI p.~106. In controversia famosa inter Maupertuis et König, a principio actionis minimae orta, sed mox ad res heterogeneas egressa, König in manibus se habere dixit autographum Leibnitianum, in quo demonstratio huius theorematis cum Euleriana prorsus conspirans contineatur. Appel au public, p.~166. Licet vero fidem huic testimonio denegare nolimus, certe Leibnitius inventum suum numquam publicavit. Conf. Hist. de l'Ac. de Prusse, A. 1750 p.~530.}. Innititur ista evolutioni potestatis $(a+1)^p$, ubi ex coefficientium forma facillime deducitur, $(a+1)^p - a^p - 1$ semper per $p$ fore divisibilem, adeoque $(a+1)^p - (a+1)$ per $p$ divisibilem fore quando $a^p - a$ per $p$ sit divisibilis.

Iam quia $1^p - 1$ semper per $p$ divisibilis est, etiam $2^p - 2$ semper erit; hinc etiam $3^p - 3$ etc. generaliterque $a^p - a$. Quodsi itaque $p$ ipsum $a$ non metitur, etiam $a^{p-1} - 1$ per $p$ divisibilis erit. Haec sufficient ad methodi indolem declarandam. Clar. Lambert similem demonstrationem tradidit in Actis Erudit. 1769 p.~109.

Quia vero evolutio potestatis binomii a theoria numerorum satis aliena esse videbatur, aliam demonstrationem ill. Euler investigavit, quae exstat Comment. nov. Petr. T.~VII p.~70 cum ea quam nos art. praec. exposuimus prorsus convenit. In sequentibus adhuc alias quaedam se nobis offerent. Hoc loco unam superaddere liceat, quae similibus principiis innititur, uti prima ill. Euleri.

\begin{satz}[Polynomii potestas.]
\setcounter{satz}{51}
\addtocounter{satz}{-1}
\textbf{51.}

Polynomii $a+b+c+\text{etc.}$ potestas $p^\text{ta}$ secundum modulum $p$ est 
\[a^p + b^p + c^p + \text{etc.},\]
siquidem $p$ est numerus primus.
\end{satz}

\textsc{Demonstr.} Constat potestatem $p^\text{tam}$ polynomii $a+b+c+\text{etc.}$ esse compositam e partibus formae $N \cdot a^\alpha b^\beta c^\gamma \text{ etc.}$, ubi $\alpha+\beta+\gamma \text{ etc.} = p$, et $N$ designat, quot modis $p$ res, quarum $\alpha$, $\beta$, $\gamma$ etc. respective sunt $= a$, $b$, $c$ etc., permutari possint. At supra art.~41 ostendimus hunc numerum semper esse per $p$ divisibilem, nisi omnes res sint aequales, i.~e. nisi aliquis numerorum $\alpha$, $\beta$, $\gamma$ etc. sit $=p$, reliqui vero $=0$. Unde sequitur, omnes ipsius $(a+b+c+\text{etc.})^p$ partes, praeter has $a^p$, $b^p$, $c^p$ etc., per $p$ divisibiles esse; quae igitur, quando de congruentia secundum modulum $p$ agitur, tuto omitti poterunt, fietque
\[(a+b+c+\text{etc.})^p \equiv a^p + b^p + c^p + \text{etc.} \pmod{p}. \quad \text{Q.~E.~D.}\]

Quodsi iam omnes quantitates $a$, $b$, $c$ etc. $=1$ ponuntur, numerusque earum $=k$, fiet $k^p \equiv k$, uti in art. praec.

\subsection*{I. Quot numeris respondeant periodi, in quibus terminorum multitudo est divisor datus numeri $p-1$.}

\setcounter{satz}{52}
\addtocounter{satz}{-1}
\textbf{52.}

Quoniam igitur alii numeri quam qui sunt divisores ipsius $p-1$, nequeunt esse exponentes potestatum infimarum, ad quas evecti numeri aliqui unitati congruunt, quaestio sese offert, num omnes ipsius $p-1$ divisores ad hoc sint idonei, atque, quando omnes numeri per $p$ non divisibiles secundum exponentem intimae suae potestatis unitati congruae classificentur, quot ad singulos exponentes sint perventuri. Ubi statim observare convenit, sufficere, si omnes numeri positivi ab $1$ usque ad $p-1$ considerentur; manifestum enim est, numeros congruos ad eandem potestatem elevari debere, quo unitati fiant congruae, adeoque numerum quemcunque ad eundem exponentem esse referendum, ad quem residuum suum minimum positivum. Quocirca in id nobis erit incumbendum ut, quomodo hoc respectu numeri $1, 2, 3, \ldots, p-1$ inter singulos factores numeri $p-1$ distribuendi sint, eruamus. Brevitatis gratia, si $d$ est unus e divisoribus numeri $p-1$ (ad quos etiam $1$ et $p-1$ referendi), per $\psi d$ designabimus multitudinem numerorum positivorum ipso $p$ minorum, quorum potestas $d^\text{ta}$ est infima unitati congrua.

\setcounter{satz}{53}
\addtocounter{satz}{-1}
\textbf{53.}

Quo facilius haec disquisitio intelligi possit, exemplum apponimus. Pro $p=19$ distribuentur numeri $1, 2, 3, \ldots, 18$ inter divisores numeri $18$ hoc modo:
\begin{center}
\begin{tabular}{c|ccccccc}
$1$ & $2$ & $3$ & $6$ & $9$ & $18$ \\
\hline
$1$ & $18$ & $7,11$ & $8,12$ & $4,5,6,9,16,17$ & $2,3,10,13,14,15$ \\
$\psi d$ & $1$ & $2$ & $2$ & $6$ & $6$
\end{tabular}
\end{center}

\setcounter{satz}{54}
\addtocounter{satz}{-1}
\textbf{54.}

II. Iam sint omnes divisores numeri $p-1$ hi: $d$, $d'$, $d''$ etc. eritque, quia omnes numeri $1, 2, 3, \ldots, p-1$ inter hos sunt distributi
\[\psi d + \psi d' + \psi d'' + \text{etc.} = p-1.\]
At in art.~40 demonstravimus esse
\[\varphi d + \varphi d' + \varphi d'' + \text{etc.} = p-1\]
atque ex art. praec. sequitur, $\psi d$ ipsi $\varphi d$ aut aequalem aut ipso minorem esse, maiorem esse non posse, similiterque de $\psi d'$ et $\varphi d'$, etc. Si itaque aliquis terminus ex his $\psi d$, $\psi d'$, $\psi d''$ etc. termino respondente ex his $\varphi d$, $\varphi d'$, $\varphi d''$ etc. esset minor (sive etiam plures), illorum summa summae horum aequalis esse non posset. Unde tandem concludimus, $\psi d$ ipsi $\varphi d$ semper esse aequalem, adeoque a magnitudine ipsius $p-1$ non pendere.

\setcounter{satz}{55}
\addtocounter{satz}{-1}
\textbf{55.}

Maximam autem attentionem meretur casus particularis propositionis praecedentis, scilicet semper dari numeros, quorum nulla potestas inferior quam $p-1^{\text{ma}}$ unitati congrua, et quidem totidem inter $1$ et $p-1$, quot infra $p-1$ sint numeri ad $p-1$ primi. Cuius theorematis demonstratio quum minime tam obvia sit quam primo aspectu videri possit, propter theorematis dignitatem liceat aliam adhuc adiicere a praecedente aliquantum diversam, quandoquidem methodorum diversitas ad res obscuriores illustrandas plurimum conferre solet. Resolvatur $p-1$ in factores suos primos fiatque $p-1 = a^\alpha b^\beta c^\gamma$ etc., designantibus $a$, $b$, $c$ etc. numeros primos inaequales. Tum theorematis demonstrationem per sequentia absolvemus.

Semper inveniri posse numerum $A$ (aut plures) ad exponentem $a^\alpha$ pertinentem similiterque numeros $B$, $C$ etc. ad exponentes $b^\beta$, $c^\gamma$ etc. respective pertinentes.

II. Productum ex omnibus numeris $A$, $B$, $C$ etc. (sive huius producti residuum minimum) ad exponentem $p-1$ pertinere.

Haec autem ita demonstramus.

I. Sit $g$ numerus aliquis ex his $1, 2, 3, \ldots, p-1$, congruentiae $x^{\frac{p-1}{a^\alpha}} \equiv 1 \pmod{p}$ non satisfaciens (omnes enim hi numeri congruentiae huic, cuius gradus $<p-1$, satisfacere nequeunt). Tum dico si potestas $\frac{p-1}{a^\alpha}$ta ipsius $g$ ponatur $=h$, hunc numerum sive eius residuum minimum ad exponentem $a^\alpha$ pertinere.

Namque patet potestatem $a^\alpha$ tam ipsius $h$ congruam fore potestati $\frac{p-1}{a^\alpha} \cdot a^\alpha = p-1$ ipsius $g$, i.~e. unitati, potestas vero $a^{\alpha-1}$ta ipsius $h$ congrua erit potestati $\frac{p-1}{a}$ ipsius $g$, i.~e. unitati erit incongrua, multoque minus potestates $a^{\alpha-2}$, $a^{\alpha-3}$ etc. ipsius $h$ unitati congruae esse possunt. At exponens infimae potestatis ipsius $h$ unitati congruae, sive exponens ad quem pertinet $h$, numerum $a^\alpha$ metiri debet (art.~48). Quare quum $a^\alpha$ per alios numeros divisibilis non sit quam per se ipsum, atque per inferiores ipsius $a$ potestates, necessario $a^\alpha$ erit exponens ad quem $h$ pertinet. Q.~E.~D. Per similem methodum demonstratur, dari numeros ad exponentes $b^\beta$, $c^\gamma$ etc. pertinentes.

II. Si supponimus, productum ex omnibus $A$, $B$, $C$ etc. non ad exponentem $p-1$, sed ad minorem $t$ pertinere, $t$ ipsum $p-1$ metietur (art.~48), sive erit $\frac{p-1}{t}$ integer unitate maior. Facile autem perspicitur, hunc quotientem vel esse unum e numeris primis $a$, $b$, $c$ etc. vel saltem per aliquem eorum divisibilem (art.~17), ex. gr. per $a$, de reliquis enim simile est ratiocinium. Metietur itaque $t$ ipsum $\frac{p-1}{a^\alpha}$; quare productum $ABC$ etc. etiam ad potestatem $\frac{p-1}{a^\alpha}$ elevatum unitati erit congruum (art.~46). Sed perspicuum est singulos $B$, $C$ etc. (excepto ipso $A$) ad potestatem $\frac{p-1}{a^\alpha}$ elevatos unitati congruos fieri, quum exponentes $b^\beta$, $c^\gamma$ etc., ad quos singuli pertinent, ipsum $\frac{p-1}{a^\alpha}$ metiantur. Hinc erit
\[A^{\frac{p-1}{a^\alpha}} B^{\frac{p-1}{a^\alpha}} C^{\frac{p-1}{a^\alpha}} \text{ etc.} \equiv A^{\frac{p-1}{a^\alpha}} \equiv 1.\]
Unde sequitur exponentem, ad quem $A$ pertinet, ipsum $\frac{p-1}{a^\alpha}$ metiri debere (art.~48), i.~e. $\frac{p-1}{a^\alpha \cdot t}$ esse integrum; at $\frac{p-1}{a^\alpha t} = \frac{1}{\text{integer}}$ integer esse nequit (art.~15). Unde tandem concludere oportet, suppositionem nostram consistere non posse, i.~e. productum $ABC$ etc. revera ad exponentem $p-1$ pertinere. Q.~E.~D.

Demonstratio posterior priori aliquantulum prolixior esse videtur, prior contra posteriori minus directa.

\setcounter{satz}{56}
\addtocounter{satz}{-1}
\textbf{56.}

Hoc theorema insigne exemplum suppeditat, quanta circumspectione in theoria numerorum saepenumero opus sit, ne, quae non sunt, pro certis assumamus. Celeb. Lambert in diss. iam supra laudata Acta Erudit. 1769 p.~127 huius propositionis mentionem facit, sed demonstrationis ne necessitatem quidem attigit. Nemo vero demonstrationem tentavit praeter summum Eulerum, Comment. nov. Acad. Petrop. T.~XVIII ad annum 1773, Demonstrationes circa residua ex divisione potestatum per numeros primos resultantia p.~85 seqq. vid. imprimis art.~37, ubi de demonstrationis necessitate fusius locutus est. At demonstratio, quam Vir sagacissimus exhibuit, duos defectus habet. Alterum quod art.~31 et sqq. tacite supponit, congruentiam $x^n \equiv 1$ (translatis ratiociniis illic adhibitis in nostra signa) revera $n$ radices diversas habere, quamquam ante nihil aliud fuerit demonstratum quam quod plures habere nequeat; alterum, quod formulam art.~34 per inductionem tantummodo deduxit.


\section*{Radices primitivae, bases, indices.}

\setcounter{satz}{57}
\addtocounter{satz}{-1}
\textbf{57.}

Numeros ad exponentem $p-1$ pertinentes \emph{radices primitivas} cum ill. Euleri vocabimus. Si igitur $a$ est radix primitiva, potestatum $a$, $a^2$, $a^3$, \ldots, $a^{p-1}$ residua minima omnia erunt diversa; unde facile deducitur, inter haec omnes numeros $1, 2, 3, \ldots, p-1$, qui totidem sunt multitudine quot illa residua minima, reperiri debere: i.~e. quemvis numerum per $p$ non divisibilem potestati alicui ipsius $a$ congruum esse. Insignis haec proprietas permagnae est utilitatis, operationesque arithmeticas, ad congruentias pertinentes, haud parum sublevare potest, simili fere modo, ut logarithmorum introductio operationes arithmeticae vulgaris. Radicem aliquam primitivam, $a$, ad lubitum pro \emph{basi} adoptabimus, ad quam omnes numeros per $p$ non divisibiles referemus, et si fuerit $a^\mu \equiv b \pmod{p}$, $\mu$ ipsius $b$ \emph{indicem} vocabimus. Ex. gr. si pro modulo $19$, radix primitiva $2$ pro basi assumitur, respondebunt

\begin{center}
\begin{tabular}{l|cccccccccccccccccc}
numeris & 1 & 2 & 3 & 4 & 5 & 6 & 7 & 8 & 9 & 10 & 11 & 12 & 13 & 14 & 15 & 16 & 17 & 18 \\
indices & 0 & 1 & 13 & 2 & 16 & 14 & 6 & 3 & 8 & 17 & 12 & 15 & 5 & 7 & 11 & 4 & 10 & 9
\end{tabular}
\end{center}

Ceterum patet, manente basi, cuique numero plures indices convenire, sed hos omnes secundum modulum $p-1$ fore congruos; quamobrem quoties de indici-bus sermo erit, qui secundum modulum $p-1$ sunt congrui, pro aequivalentibus habebuntur, simili modo uti numeri ipsi, quando secundum modulum $p$ sunt congrui, tamquam aequivalentes spectantur.

\subsection*{Algorithmus indicum.}

\setcounter{satz}{58}
\addtocounter{satz}{-1}
\textbf{58.}

Theoremata ad indices pertinentia prorsus analoga sunt iis, quae ad logarithmos spectant.

Index producti e quotcunque factoribus conflati congruus est summae indicum singulorum factorum secundum modulum $p-1$.

Index potestatis numeri alicuius congruus est producto ex indice numeri dati in exponentem potestatis, secundum mod.~$p-1$.

Demonstrationes propter facilitatem omittimus.

Hinc perspicitur, si tabulam construere velimus, ex qua omnium numerorum indices pro modulis diversis desumi possint, ex hac tum omnes numeros modulo maiores, tum omnes compositos omitti posse. Specimen huius modi tabulae ad calcem operis huius adiectum est, Tab.~I, ubi in prima columna verticali positi sunt numeri primi primorumque potestates a $3$ usque ad $97$, qui tamquam moduli sunt spectandi; iuxta hos singulos numeri pro basi assumti; tum sequuntur indices numerorum primorum successivorum, quorum quini semper per parvulum intervallum sunt disiuncti, eodemque ordine supra dispositi sunt numeri primi; ita ut quis index numero primo dato secundum modulum datum respondeat, facile tutoque inveniri possit.

Ita ex. gr. si $p=67$ index numeri $60$, assumto $12$ pro basi erit
\[= \text{Ind. } 2 + \text{Ind. } 3 + \text{Ind. } 5 \pmod{66} = 58 + 9 + 39 = 40.\]

\setcounter{satz}{59}
\addtocounter{satz}{-1}
\textbf{59.}

Index valoris cuiuscunque expressionis $\frac{a}{b} \pmod{p}$ (art.~31) congruus est secundum modulum $p-1$ differentiae indicum numeratoris $a$ et denominatoris $b$, siquidem numeri $a$, $b$ per $p$ non sunt divisibiles.

Sit enim valor quicunque $c$: eritque $bc \equiv a \pmod{p}$; hinc
\[\text{Ind. } b + \text{Ind. } c \equiv \text{Ind. } a \pmod{p-1}\]
adeoque
\[\text{Ind. } c \equiv \text{Ind. } a - \text{Ind. } b.\]

Si itaque tabula habetur, ex qua index cuique numero respondens pro quovis modulo primo, aliaque ex qua numerus ad indicem datum pertinens derivari possit, omnes congruentiae primi gradus facillimo negotio solvi poterunt, quoniam omnes reduci possunt ad tales, quarum modulus est numerus primus (art.~30).

E.~g. proposita congruentia $29x + 7 \equiv 0 \pmod{47}$, erit $x \equiv \frac{-7}{29} \pmod{47}$. Hinc
\[\text{Ind. } x \equiv \text{Ind. }(-7) - \text{Ind. } 29 = \text{Ind. } 40 - \text{Ind. } 29 = 15 - 43 = 18 \pmod{46}.\]
At numerus cuius index $18$ invenitur $3$. Quare $x \equiv 3 \pmod{47}$.

Tabulam secundam quidem non adiecimus: at huius vice alia defungi poterit, uti Sect.~VI ostendemus.

\subsection*{De radicibus congruentiae $x^n \equiv A$.}

\setcounter{satz}{60}
\addtocounter{satz}{-1}
\textbf{60.}

Simili modo ut art.~31 radices congruentiarum primi gradus designavimus, in sequentibus etiam congruentiarum purarum altiorum graduum radices per signum exhibebimus. Ut scilicet $\sqrt[n]{A}$ nihil aliud significat quam radicem aequationis $x^n = A$, ita apposito modulo per $\sqrt[n]{A} \pmod{p}$ denotabitur radix quaecunque congruentiae $x^n \equiv A \pmod{p}$. Hanc expressionem $\sqrt[n]{A} \pmod{p}$ tot valores habere dicemus, quot habet secundum $p$ incongruos, omnes enim qui secundum $p$ sunt congrui tamquam aequivalentes spectandi (art.~26). Ceterum patet, si $A$, $B$ secundum $p$ fuerint congrui, expressiones $\sqrt[n]{A}$, $\sqrt[n]{B} \pmod{p}$ aequivalentes fore.

Iam si ponitur $\sqrt[n]{A} \equiv x \pmod{p}$, erit $n \cdot \text{Ind. } x \equiv \text{Ind. } A \pmod{p-1}$. Ex hac congruentia deducuntur ad praecepta sectionis praec. valores ipsius Ind.~$x$ atque ex his valores respondentes ipsius $x$. Facile vero perspicitur, $x$ habere totidem valores, quot radices congruentia $n \cdot \text{Ind. } x \equiv \text{Ind. } A \pmod{p-1}$. Manifesto igitur $\sqrt[n]{A}$ unum tantummodo valorem habebit, quando $n$ ad $p-1$ est primus; quando vero numeri $n$, $p-1$ divisorem communem habent $\delta$, atque hic est maximus, Ind.~$x$ habebit $\delta$ valores incongruos secundum $p-1$, adeoque $\sqrt[n]{A}$ totidem valores incongruos secundum $p$, siquidem Ind.~$A$ per $\delta$ est divisibilis. Qua conditione deficiente $\sqrt[n]{A}$ nullum valorem realem habebit.

\textsc{Exemplum.} Quaeruntur valores expressionis $\sqrt[3]{11} \pmod{19}$. Solvi itaque debet congruentia $3 \cdot \text{Ind. } x \equiv \text{Ind. } 11 = 6 \pmod{18}$, invenienturque tres valores ipsius Ind.~$x$, $2$, $8$, $14$ (mod.~$18$). His vero respondent valores ipsius $x$, $6$, $9$, $4$.


\setcounter{satz}{61}
\addtocounter{satz}{-1}
\textbf{61.}

Quantumvis expedita sit methodus haec, quando tabulae necessariae adsunt, debemus tamen non oblivisci, indirectam eam esse. Operae igitur pretium erit inquirere quantum methodi directae polleant: trademusque hic ea quae ex praecedentibus hauriri possunt; alia, quae considerationes reconditiores postulant, ad sectionem VIII reservantes. Initium facimus a casu simplicissimo, ubi $A=1$, sive ubi radices congruentiae $x^n \equiv 1 \pmod{p}$ quaeruntur. Hic itaque, assumta radice quacunque primitiva pro basi, debet esse $n \cdot \text{Ind. } x \equiv 0 \pmod{p-1}$. Quae congruentia, quando $n$ ad $p-1$ est primus, unam tantummodo radicem habebit, scilicet Ind.~$x \equiv 0$ (mod.~$p-1$): quare in hocce casu $\sqrt[n]{1} \pmod{p}$ unicum valorem habet, scilicet $1$. Quando autem numeri $n$, $p-1$ habent divisorem communem (maximum) $\delta$, congruentiae $n \cdot \text{Ind. } x \equiv 0 \pmod{p-1}$ solutio completa erit Ind.~$x \equiv 0 \pmod{\frac{p-1}{\delta}}$ (V. art.~29), i.e. Ind.~$x$ secundum modulum $p-1$ alicui ex his numeris
\[0, \quad \frac{p-1}{\delta}, \quad \frac{2(p-1)}{\delta}, \quad \frac{3(p-1)}{\delta}, \quad \ldots, \quad \frac{(\delta-1)(p-1)}{\delta}\]
congruus esse debebit, sive $x$ in hocce casu $\delta$ valores diversos (secundum modulum $p$ incongruos) habebit. Hinc perspicitur, expressionem $\sqrt[n]{1}$ etiam $\delta$ valores diversos habere, quorum indices cum ante allatis prorsus conveniant. Quocirca expressio $\sqrt[n]{1} \pmod{p}$ huic $\sqrt[\delta]{1} \pmod{p}$ omnino aequivalet i.~e. congruentia $x^n \equiv 1 \pmod{p}$ easdem radices habet quam haec, $x^\delta \equiv 1 \pmod{p}$. Prior autem inferioris erit gradus, siquidem $\delta$ et $n$ sunt inaequales.

Ex. $\sqrt[15]{1} \pmod{19}$ tres habet valores, quia $3$ maxima numerorum $15$, $18$ mensura communis, hique simul erunt valores expressionis $\sqrt[3]{1} \pmod{19}$. Sunt autem hi $1$, $7$, $11$.

\setcounter{satz}{62}
\addtocounter{satz}{-1}
\textbf{62.}

Per hanc igitur reductionem id lucramur, ut alias congruentias formae $x^n \equiv 1$ solvere non sit opus, quam ubi $n$ numeri $p-1$ est divisor. Infra vero ostendemus, congruentias huius formae semper ulterius adhuc deprimi posse, licet praecedentia ad hoc non sufficiant. Unum tamen casum iam hic absolvere possumus, scilicet ubi $n=2$. Manifeste enim valores expressionis $\sqrt[2]{1}$ erunt $+1$ et $-1$, quia plures quam duos habere nequit, hique $+1$ et $-1$ semper sunt incongrui, nisi modulus sit $=2$, in quo casu $\sqrt[2]{1}$ unum tantum valorem habere posse per se clarum. Hinc sequitur, $+1$ et $-1$ etiam fore valores expressionis $\sqrt[n]{1}$, quando $n$ ad $\frac{p-1}{2}$ sit primus. Hoc semper eveniet quoties modulus est eius indolis, ut $\frac{p-1}{2}$ numerus absolute primus (nisi forte $p=1=2m$, in quo casu omnes numeri $1, 2, \ldots, p-1$ sunt radices) ex. gr. quando $p=3$, $5$, $7$, $11$, $23$, $47$, $59$, $83$, $107$ etc.

Tamquam corollarium hic annotetur, indicem ipsius $-1$ semper esse $\frac{p-1}{2} \pmod{p-1}$, quaecunque radix primitiva pro basi accipiatur. Namque $2 \cdot \text{Ind.}(-1) \equiv 0 \pmod{p-1}$. Quare Ind.~$(-1)$ erit vel $0$, vel $\frac{p-1}{2} \pmod{p-1}$: vero semper index ipsius $+1$, atque $+1$ et $-1$ semper indices diversos habere debent (praeter casum $p=2$, ad quem hie respicere operae non est pretium).

\setcounter{satz}{63}
\addtocounter{satz}{-1}
\textbf{63.}

Ostendimus art.~60, expressionem $\sqrt[n]{A} \pmod{p}$ habere $\delta$ valores diversos, aut omnino nullum, si fuerit $\delta$ divisor communis maximus numerorum $n$, $p-1$. Iam uti modo docuimus $\sqrt[n]{A}$ et $\sqrt[\delta]{A}$ aequivalentes esse, si fuerit $A=1$, generalius probabimus, expressionem $\sqrt[n]{A}$ semper ad aliam $\sqrt[n]{B}$ reduci posse, cui aequivaleat. Illius enim valore quocunque denotato per $x$ erit $x^n \equiv A$; iam sit $t$ valor quicunque expressionis $\sqrt[n]{1} \pmod{p-1}$, quam valores reales habere ex art.~31 perspicuum; eritque $x^{tn} \equiv A^t$; at $x^{tn} \equiv x^n$ propter $tn \equiv n \pmod{p-1}$. Quare $x^n \equiv A^t$ adeoque quicunque ipsius $\sqrt[n]{A}$ valor erit etiam valor ipsius $\sqrt[n]{A^t}$. Quoties igitur $\sqrt[n]{A}$ valores reales habet, expressioni $\sqrt[n]{A^t}$ prorsus aequivalens erit, quoniam illa neque alios habet quam haec neque pauciores, licet quando $\sqrt[n]{A}$ nullum valorem realem habet, fieri tamen possit, ut $\sqrt[n]{A^t}$ valores reales habeat.

\textsc{Ex.} Si valores expressionis $\sqrt[21]{2} \pmod{31}$ quaeruntur, erit numerorum $21$ et $30$ divisor communis maximus $3$, expressionisque $\sqrt[3]{2^7} \pmod{30}$ valor aliquis $3$, quare si $\sqrt[21]{2}$ valores reales habet, huic expressioni $\sqrt[3]{128}$ sive $\sqrt[3]{8}$ aequivalebit, invenieturque revera, posterioris expressionis valores, qui sunt $2$, $10$, $19$, etiam priori satisfacere.

\setcounter{satz}{64}
\addtocounter{satz}{-1}
\textbf{64.}

Ne autem hanc operationem incassum suscepisse periclitemur regulam investigare oportet, per quam statim diiudicari possit, utrum $\sqrt[n]{A}$ valores reales admittat necne. Quodsi tabula indicum habetur, res in promtu est; namque ex art.~60 manifestum est, valores reales dari, si ipsius $A$ index, radice quacunque primitiva pro basi accepta, per $\delta$ sit divisibilis, sin vero minus, non dari. Attamen hoc etiam absque tali tabula inveniri potest. Posito enim indice ipsius $A = A'$, si hic fuerit per $\delta$ divisibilis, erit $\frac{n \cdot A'}{\delta}$ per $p-1$ divisibilis et vice versa.

Atqui numeri $A' \cdot \frac{p-1}{\delta}$ index erit $\frac{n \cdot A'}{\delta}$, unitati congruus erit, sin minus incongruus. Quare si $A^{\frac{n \cdot (p-1)}{2\delta}} \equiv 1 \pmod{p}$, $\sqrt[n]{A} \pmod{p}$ habet valores reales, sin minus non habet. Ita in exemplo art. praec. $2^{10} = 1024 \equiv 1 \pmod{31}$, unde concluditur $\sqrt[21]{2} \pmod{31}$ valores reales habere. Similiter certiores hinc fimus $\sqrt[\frac{p-1}{2}]{-1} \pmod{p}$ semper valores binos reales habere, quando $p$ sit formae $4m+1$, nullum vero, quando $p$ sit formae $4m+3$; propter $(-1)^{\frac{p-1}{2}} \equiv +1$ et $(-1)^{\frac{p-1}{2}} \equiv -1$. Elegans hoc theorema, quod vulgo ita profertur: Si $p$ est numerus primus formae $4m+1$, inveniri potest quadratum $aa$, ita ut $aa+1$ per $p$ fiat divisibilis; si vero $p$ est formae $4m-1$, tale quadratum non datur, hoc modo demonstratum est ab ill. Eulero, Comm. nov. Acad. Petrop. T.~XVIII p.~112 ad annum 1773. Demonstrationem aliam iam multo ante dederat, Comm. nov. T.~V p.~5, qui prodiit a.~1760. In dissert. priori, Comm. nov. T.~IV p.~25, rem nondum perfecerat. Postea etiam ill. La Grange theorematis demonstrationem tradidit Nouv. Mém. de l'Ac. de Berlin A.~1775 p.~342. Aliam adhuc demonstrationem in sectione sequenti, ubi proprie de hoc argumento agendum erit, dabimus.

\setcounter{satz}{65}
\addtocounter{satz}{-1}
\textbf{65.}

Postquam omnes expressiones $\sqrt[n]{A} \pmod{p}$ ad tales reducere docuimus, ubi $n$ divisor numeri $p-1$, criteriumque nacti sumus, utrum valores reales admittat necne, tales expressiones $\sqrt[n]{A} \pmod{p}$, ubi $n$ ipsius $p-1$ est divisor, accuratius considerabimus. Primo ostendemus, quam relationem valores singuli expressionis inter se habeant, tum artificia quaedam trademus, quorum auxilio unus valor expressionis saepenumero inveniri possit.

\textsc{Primo.} Quando $A \equiv 1$ atque $r$ aliquis ex $n$ valoribus expressionis $\sqrt[n]{1} \pmod{p}$, sive $r^n \equiv 1 \pmod{p}$, omnes etiam ipsius $r$ potestates erunt valores istius expressionis; horum autem totidem erunt diversi, quot unitates habet exponens, ad quem $r$ pertinet (art.~48). Quodsi igitur $r$ est valor ad exponentem $n$ pertinens, potestates ipsius $r$ hae $r$, $r^2$, $r^3$, \ldots, $r^n$ (ubi loco ultimae unitas substitui potest) omnes expressionis $\sqrt[n]{1} \pmod{p}$ valores involvent. Qualia autem subsidia exstent ad tales valores inveniendos, qui ad exponentem $n$ pertineant, in Sect.~VIII fusius explicabimus.

\textsc{Secundo.} Quando $A$ unitati est incongruus, unusque valor expressionis $\sqrt[n]{A} \pmod{p}$ notus, qui sit $z$, reliqui hoc modo inde deducuntur. Sint $t$ valores expressionis $\sqrt[n]{1}$ hi $r$, $r^2$, $r^3$, \ldots, $r^{n-1}$ (uti modo ostendimus), eruntque omnes expr.~$\sqrt[n]{A}$ valores hi $z$, $zr$, $zr^2$, \ldots, $zr^{n-1}$.


\setcounter{satz}{66}
\addtocounter{satz}{-1}
\textbf{66.}

Secundum quod nobis proposueramus fuit docere, in quo casu unus expressionis $\sqrt[n]{A} \pmod{p}$ valor (ubi $n$ supponitur esse divisor ipsius $p-1$) directe inveniri possit. Hoc evenit, quando aliquis valor potestati alicui ipsius $A$ congruus evadit, qui casus quum haud raro occurrat, aliquantum huic rei immorari non superfluum erit. Sit talis valor, si quis datur, $z$, sive $z^n \equiv A^k$ et $A \equiv z^m \pmod{p}$. Hinc colligitur $A \equiv A^{kn}$, quare si numerus $k$ habetur, ita ut sit $A \equiv A^{kn}$, $A^k$ erit valor quaesitus. At huic conditioni aequivalet ista ut sit $kn \equiv 1 \pmod{t}$, designante $t$ exponentem, ad quem pertinet $A$ (artt.~46,~48). Ut vero haec congruentia possibilis sit, requiritur, ut sit $k$ ad $t$ primus. Hoc in casu erit $k \equiv \frac{1}{n} \pmod{t}$; si vero $t$ et $n$ divisorem communem habent, nullus valor $z$ potestati ipsius $A$ congruus esse potest.

\setcounter{satz}{67}
\addtocounter{satz}{-1}
\textbf{67.}

Quum autem ad hanc solutionem ipsum $t$ novisse oporteat, videamus quomodo procedere possimus, si hunc numerum ignoremus. Primo facile intelligitur, $t$ ipsum $\frac{p-1}{\delta}$ metiri debere, siquidem $\sqrt[n]{A} \pmod{p}$ valores reales habeat, uti hic semper supponimus. Sit enim quicunque valor $y$, eritque tum $y^{p-1} \equiv 1$, tum $y^n \equiv A \pmod{p}$; quare elevando partes posterioris congruentiae ad potestatem $\frac{p-1}{\delta}$, fiet $A^{\frac{p-1}{\delta}} \equiv 1$; adeoque $\frac{p-1}{\delta}$ per $t$ divisibilis (art.~48). Iam si $\frac{p-1}{\delta}$ ad $n$ est primus congruentia art. praec. $kn \equiv 1$ etiam secundum modulum $\frac{p-1}{\delta}$ solvi poterit manifestoque valor ipsius $k$ congruentiae secundum modulum hunc satisfaciens eidem etiam secundum modulum $t$, qui ipsum $\frac{p-1}{\delta}$ metitur, satisfaciet (art.~5). Tum igitur quod quaerebatur, inventum.

Si vero $\frac{p-1}{\delta}$ ad $n$ non est primus, omnes ipsius $\frac{p-1}{\delta}$ factores primi, qui simul ipsum $t$ metiuntur, ex $\frac{p-1}{\delta}$ eiiciantur. Hinc nanciscemur numerum $\frac{p-1}{\delta q}$, ad $n$ primum, designante $q$ productum ex omnibus illis factoribus primis, quos eiicimus. Quodsi iam conditio ad quam in artic. praec. pervenimus ut $t$ ad $n$ sit primus, locum habet, $t$ etiam ad $q$ erit primus adeoque etiam ipsum $\frac{p-1}{\delta q}$ metietur. Quare si congruentia $kn \equiv 1 \pmod{\frac{p-1}{\delta q}}$ solvitur (quod fieri potest quia $n$ ad $\frac{p-1}{\delta q}$ primus), valor ipsius $k$ etiam secundum modulum $t$ congruentiae satisfaciet, id quod quaerebatur. Totum hoc artificium in eo versatur, ut numerus eruatur, qui ipsius $t$, quem ignoramus, vice fungi possit. Attamen probe meminisse oportet, nos, quando $\frac{p-1}{\delta}$ ad $n$ non est primus, supposuisse conditionem art. praec. locum habere, quae si deficit omnes conclusiones erroneae erunt; atque si regulas datas temere sequendo pro $z$ valor invenitur, cuius potestas $A^k$ ipsi $A$ non sit congrua, indicio hoc est, conditionem deficere adeoque methodum hanc omnino adhiberi non posse.

\setcounter{satz}{68}
\addtocounter{satz}{-1}
\textbf{68.}

Sed in hocce etiam casu saepe prodesse potest, hunc laborem suscepisse; operaeque pretium est, quomodo hic valor falsus ad veros sese habeat, investigare. Supponamus itaque, numeros $k$, $z$ rite esse determinatos sed $z^n$ non esse $A \pmod{p}$. Tum si modo valores expressionis $\sqrt[n]{\frac{A}{z^n}} \pmod{p}$ determinari possint, hos singulos per $z$ multiplicando valores ipsius $\sqrt[n]{A}$ obtinebimus. Si enim $v$ est valor aliquis ipsius $\sqrt[n]{\frac{A}{z^n}}$, erit $(vz)^n \equiv A$. Sed expressio $\sqrt[n]{\frac{A}{z^n}}$ eatenus hac $\sqrt[n]{A}$ simplicior, quod $\frac{A}{z^n} \pmod{p}$ ad exponentem minorem plerumque pertinet quam $A$. Scilicet si numerorum $t$, $q$ divisor communis maximus est $d$, $\frac{A}{z^n} \pmod{p}$ ad exponentem $d$ pertinebit, id quod ita demonstratur. Substitute pro $z$ valore, fit $\frac{A}{z^n} \equiv A^{1-kn} \pmod{p}$. At $kn-1$ per $\frac{p-1}{\delta q}$ divisibilis (art. praec), $\frac{p-1}{\delta}$ vero per $t$ (ibid.) sive $\frac{t}{d}$. Atqui $\frac{t}{d}$ ad $\frac{q}{d}$ est primus (hyp.), quare etiam $kn-1$ per $\frac{t}{d}$ sive $kn-1$ per $\frac{t}{d}$, adeoque etiam $(kn-1)d$ per $t$ erit divisibilis. Hinc $A^{(kn-1)d} \equiv 1 \pmod{p}$. Unde facile deducitur, $\frac{A}{z^n}$ ad potestatem $d^\text{tam}$ evectum unitati congruum fieri. Quod vero $\frac{A}{z^n}$ ad exponentem minorem quam $d$ pertinere non possit, facile quidem demonstrari potest, sed quoniam ad finem nostrum non requiritur, huic rei non immoramur. Certi igitur esse possumus, $\frac{A}{z^n} \pmod{p}$ semper ad minorem exponentem pertinere quam $A$, unico excepto casu, scilicet quando $t$ ipsum $q$ metitur, adeoque $d=t$.

\setcounter{satz}{69}
\addtocounter{satz}{-1}
\textbf{69.}

Revertimur nunc ad radices, quas diximus primitivas. Ostendimus, radice primitiva quacunque pro basi assumta omnes numeros, quorum indices ad $p-1$ primi, etiam fore radices primitivas, nullosque praeter hos: unde simul radicum primitivarum multitudo sponte innotescit. V. art.~53. Quamnam autem radicem primitivam pro basi adoptare velimus, in genere arbitrio nostro relinquitur; unde intelligitur, etiam hic, ut in calculo logarithmico, plura quasi systemata dari posse\footnote{In eo autem differunt, quod in logarithmis systematum numerus est infinitus, hic vero tantus, quantus numerus radicum primitivarum. Manifeste enim bases congruae idem systema generant.}, quae quo vinculo connexa sint videamus. Sint $a$, $b$ duae radices primitivae, aliusque numerus $m$, atque, quando $a$ pro basi assumitur, index numeri $b \equiv \mu$, numeri $m$ vero index $\equiv \alpha \pmod{p-1}$; quando autem $b$ pro basi assumitur, index numeri $a \equiv \alpha'$, numeri $m$ vero $\equiv \nu \pmod{p-1}$. Tum erit $a^{\alpha'} \equiv b$, quare $a^{\alpha'\mu} \equiv b^\mu \equiv a \pmod{p}$ (hyp.), hinc $a^{\alpha'\mu-1} \equiv 1 \pmod{p-1}$. Per simile ratiocinium invenitur $\nu \equiv \alpha \mu'$, atque $\mu \mu' \equiv 1 \pmod{p-1}$. Si igitur tabella indicum pro basi $a$ constructa habetur, facile in aliam converti potest, ubi $b$ basis. Si enim pro basi $a$ ipsius $b$ index est $=\varepsilon$, pro basi $b$ ipsius $a$ index erit $\frac{1}{\varepsilon} \pmod{p-1}$, multiplicandoque per hunc numerum omnes tabellae indices, habebuntur omnes indices pro basi $b$.

\setcounter{satz}{70}
\addtocounter{satz}{-1}
\textbf{70.}

Quamvis autem plures indices numero dato contingere possint, aliis aliisque radicibus primitivis pro basi acceptis, omnes tamen in eo convenient, quod omnes eundem divisorem maximum cum $p-1$ communem habebunt. Si enim pro basi $a$, index numeri dati est $m$, pro basi $b$ vero $n$, atque divisores maximi his cum $p-1$ communes $\mu$, $\nu$ supponuntur esse inaequales, alter erit maior, ex. gr. $\mu > \nu$, adeoque $\mu$ ipsum $n$ non metietur. At designato indice ipsius $a$, quando $b$ pro basi assumitur, per $\alpha$, erit (art. praec.) $n \equiv \alpha m \pmod{p-1}$ adeoque $\mu$ etiam ipsum $n$ metietur. Q.~E.~A.

Hunc divisorem maximum indicibus numeri dati, ipsique $p-1$ communem, a basi non pendere, etiam inde perspicuum, quod aequalis est ipsi $\frac{p-1}{t}$, designante $t$ exponentem ad quem numerus, de cuius indicibus agitur, pertinet. Si enim index pro basi quacunque est $k$, erit $t$ minimus numerus per quem $k$ multiplicatus ipsius $p-1$ multiplum evadit (excepta cifra) vid. artt.~48,~58, sive minimus valor expressionis $\frac{0}{k} \pmod{p-1}$ praeter cifram; hunc autem aequalem esse divisori maximo communi numerorum $k$ et $p-1$, ex art.~29 nullo negotio derivatur.

\setcounter{satz}{71}
\addtocounter{satz}{-1}
\textbf{71.}

Porro facile demonstratur, basin ita semper accipere licere, ut numerus ad exponentem $t$ pertinens indicem quemlibet datum nanciscatur, cuius quidem maximus divisor cum $p-1$ communis $=d$. Designemus hunc brevitatis gratia per $d$, sitque index propositus $= dm$, numerique propositi, quando quaelibet radix primitiva $a$ pro basi accipitur, index $\equiv dn$ eruntque $m$, $n$ ad $\frac{p-1}{d}$ sive ad $t$ primi. Tum si $s$ est valor expressionis $\frac{m}{n} \pmod{p-1}$, simulque ad $p-1$ primus, erit $a^s$ radix primitiva, qua pro basi accepta numerus propositus indicem $dm$ adipiscetur (erit enim $a^{sdn} \equiv a^{dm} \equiv$ numero proposito), id quod desiderabatur. Sed expressionem $\frac{m}{n} \pmod{p-1}$ valores ad $p-1$ primos admittere, ita probatur. Aequivalet illa expressio huic: $\frac{m}{n} \pmod{\frac{p-1}{d}}$ sive $\frac{m}{n} \pmod{t}$ vid. art.~31,~2, eruntque omnes eius valores ad $t$ primi; si enim aliquis valor e divisorem cum $t$ communem haberet, hic divisor etiam ipsum $m$ metiri deberet, adeoque etiam ipsum $n$, cui $m$ secundum $t$ congruus, contra hypoth., ex qua $m$ ad $t$ primus.

Quando igitur omnes divisores primi ipsius $p-1$ etiam ipsum $t$ metiuntur, omnes expr. $\frac{m}{n} \pmod{t}$ valores ad $p-1$ primi erunt multitudoque eorum $=\varphi t$; quando autem $p-1$ alios adhuc divisores primos, $f$, $g$, $h$ etc. implicat, ipsum $t$ non metientes, ponatur valor quicunque expr.~$\frac{m}{n} \pmod{t} = e$. Tum autem quia omnes $t$, $f$, $g$, $h$ etc. inter se primi, inveniri potest numerus $\varepsilon$, qui secundum $t$ ipsi $e$, secundum $f$, $g$, $h$ etc. vero numeris quibuscunque ad hos respective primis fiat congruus (art.~32). Talis itaque numerus per nullum factorem primum ipsius $p-1$ divisibilis adeoque ad $p-1$ primus erit, uti desiderabatur. Tandem haud difficile ex combinationum theoria deducitur, talium valorum multitudinem fore $= \frac{\varphi t \cdot \varphi f \cdot \varphi g \ldots}{\varphi(t \cdot f \cdot g \ldots)}$ sed ne digressio haec in nimiam molem excrescat, demonstrationem, quum ad institutum nostrum non sit adeo necessaria, omittimus.


\subsection*{Bases usibus peculiaribus accommodatae.}

\setcounter{satz}{72}
\addtocounter{satz}{-1}
\textbf{72.}

Quamvis in genere prorsus arbitrarium sit quaenam radix primitiva pro basi adoptetur, interdum tamen bases aliae prae aliis commoda quaedam peculiaria praebere possunt. In tabula I semper numerum $10$ pro basi assumsimus, quando fuit radix primitiva; alioquin basin ita semper determinavimus, ut numeri $10$ index evaserit quam minimus, i.~e. $=\frac{d}{2}$, denotante $t$ exponentem, ad quem $10$ pertinuit. Quid vero hinc lucremur, in Sect.~VI ostendemus, ubi eadem tabula ad alios adhuc usus adhibebitur. Sed quoniam etiam hic aliquid arbitrarii remanere potest, ut ex art. praec. apparet: ut aliquid certi statueremus, ex omnibus radicibus primitivis quaesitum praestantibus minimam semper pro basi elegimus. Ita pro $p=41$, ubi $t=8$ atque $d=9$, $a^d$ habet $\varphi 9$ i.~e. $6$ valores, qui sunt $5$, $14$, $20$, $28$, $39$, $40$. Assumsimus itaque minimum $5$ pro basi.

\subsection*{Methodus radices primitivas assignandi.}

\setcounter{satz}{73}
\addtocounter{satz}{-1}
\textbf{73.}

Methodi radices primitivas inveniendi maximam partem tentando innituntur. Si quis ea quae art.~55 docuimus cum iis quae infra de solutione congruentiae $x^n \equiv 1$ trademus confert, omnia fere quae per methodos directas effici possunt, habebit. Ill. Euler confitetur, Opusc. Analyt. T.~I.~p.~152, maxime difficile videri, hos numeros assignare, eorumque indolem ad profundissima numerorum mysteria esse referendam. At tentando satis expedite sequenti modo determinari possunt. Exercitatus operationis prolixitati per multifaria artificia particularia succurrere sciet: haec vero per usum multo citius quam per praecepta ediscuntur.

$1^\circ$. Assumatur ad libitum numerus ad $p$ (ita semper modulum designamus) primus, $a$, (plerumque ad calculi brevitatem conducit, si quam minimum accipimus, ex. gr. numerum $2$) determineturque eius periodus (art.~46), i.~e. residua minima ipsius potestatum, donec ad potestatem $a^t$ perveniatur, cuius residuum minimum sit $=1$\footnote{Quisquis sponte perspiciet, non opus esse has potestates ipsas novisse, quum cuiusvis residuum minimum facile ex residuo minimo potestatis praecedentis obtineri possit.}. Iam si fuerit $t=p-1$, $a$ est radix primitiva.

$2^\circ$. Si vero $t < p-1$, accipiatur alius numerus $b$ in periodo ipsius $a$ non contentus, investigeturque simili modo huius periodus. Designato exponente ad quem $b$ pertinet per $u$, facile perspicitur $u$ neque ipsi $t$ aequalem neque ipsius partem aliquotam esse posse, in utroque enim casu fieret $b^t \equiv 1$, quod esse nequit, quum periodus ipsius $a$ omnes numeros amplectatur, quorum potestas exponentis $t$ unitati congrua (art.~53). Quodsi $u$ fuerit $=p-1$, erit $b$ radix primitiva; si vero $u$ non quidem $=p-1$, sed tamen multiplum ipsius $t$, id lucrati sumus, ut numerus constet ad exponentem maiorem pertinens, adeoque scopo nostro, qui est invenire numerum ad exponentem maximum pertinentem, propiores iam simus. Si vero $u$ neque $=p-1$, neque ipsius $t$ multiplum, tamen numerum invenire possumus ad exponentem ipsis $t$, $u$ maiorem pertinentem, nempe ad exponentem minimo dividuo communi numerorum $t$, $u$ aequalem. Sit hic $=y$, resolvaturque $y$ ita in duos factores inter se primos, $m$, $n$, ut alter ipsum $t$, alter ipsum $u$ metiatur\footnote{Quomodo hoc fieri possit, ex art.~18 haud difficulter derivatur. Resolvatur $y$ in factores tales, qui sint aut numeri primi diversi aut numerorum primorum diversorum potestates. Horum quisque alterutrum numerorum $t$, $u$ metietur (sive etiam utrumque). Adscribantur singuli aut numero $t$ aut numero $u$, prout illum aut hunc metiuntur: quando aliquis utrumque metitur, arbitrarium est, cui adscribatur: productum ex iis qui ipsi $t$ adscripti sunt, sit $=m$, productum e reliquis $=n$, facileque perspicietur $m$ ipsum $t$, $n$ ipsum $u$ metiri, atque esse $mn=y$.}. Tum fiat potestas $\frac{t}{m}$ta ipsius $a$, $=A$, potestas $\frac{u}{n}$ta ipsius $b$, $=B \pmod{p}$, eritque productum $AB$ numerus ad exponentem $y$ pertinens: facile enim intelligitur, $A$ ad exponentem $m$, $B$ ad exponentem $n$ pertinere; adeoque productum $AB$ ad $mn$ pertinebit, quia $m$, $n$ inter se sunt primi, id quod prorsus eodem modo uti in art.~55,~II processimus probari poterit.

Iam si $y=p-1$, $AB$ erit radix primitiva; sin minus, simili modo ut antea alius numerus adhibendus erit, in periodo ipsius $AB$ non occurrens; eritque hic aut radix primitiva aut pertinebit ad exponentem ipso $y$ maiorem aut certe ipsius auxilio (uti ante) numerus ad exponentem ipso $y$ maiorem pertinens inveniri poterit. Quum igitur numeri qui per repetitionem huius operationis prodeunt, ad exponentes continuo crescentes pertineant, manifestum est tandem numerum inventum iri, qui ad exponentem maximum pertineat, i.~e. radicem primitivam. Q.~E.~F.

\setcounter{satz}{74}
\addtocounter{satz}{-1}
\textbf{74.}

Per exemplum praecepta haec clariora fient. Sit $p=73$, pro quo radix primitiva quaeratur. Tentemus primo numerum $2$, cuius periodus prodit haec:

$1$, $2$, $4$, $8$, $16$, $32$, $64$, $55$, $37$, $1$ etc.

$0$, $1$, $2$, $3$, $4$, $5$, $6$, $7$, $8$, $9$ etc.

Quum igitur iam potestas exponentis $9$ unitati congrua fiat, $2$ non est radix primitiva. Tentetur alius numerus in periodo ipsius $2$ non occurrens ex. gr. $3$, cuius periodus est haec

$1$, $3$, $9$, $27$, $8$, $24$, $72$, $70$, $64$, $46$, $65$, $49$, $1$ etc.

$0$, $1$, $2$, $3$, $4$, $5$, $6$, $7$, $8$, $9$, $10$, $11$, $12$ etc.

Quare neque $3$ est radix primitiva. Exponentium autem ad quos $2$, $3$ pertinent, (i.e. numerorum $9$, $12$) dividuus communis minimus est $36$, qui in factores $9$ et $4$ ad praecepta art. praec. resolvitur. Evehendus itaque $2$ ad potestatem exponentis $\frac{9}{9}$, i.~e. numerus $2$ ipse retinendus; $3$ autem ad potestatem exponentis $3$; productum ex his est $54$, quod itaque ad exponentem $36$ pertinebit. Si denique ipsius $54$ periodus computatur numerusque in hac non contentus ex. gr. $5$ denuo tentatur, hunc esse radicem primitivam, reperietur.

\subsection*{Theoremata varia de periodis et radicibus primitivis.}

\setcounter{satz}{75}
\addtocounter{satz}{-1}
\textbf{75.}

Antequam hoc argumentum deseramus, quasdam trademus, propositiones quae ob simplicitatem suam attentione haud indignae videntur.

Productum ex omnibus terminis periodi numeri cuiusvis est $\equiv -1$, quando ipsorum multitudo, sive exponens ad quem numerus pertinet, est impar, et $\equiv +1$, quando hic exponens est par.

\textsc{Ex.} Pro modulo $13$ periodus numeri $5$ constat ex his terminis $1$, $5$, $12$, $8$ quorum productum $480 \equiv -1 \pmod{13}$.

Secundum eundem modulum periodus numeri $3$ constat e terminis $1$, $3$, $9$ quorum productum $27 \equiv +1 \pmod{13}$.

\textsc{Demonstr.} Sit exponens, ad quem numerus pertinet, $=t$, atque index numeri, $=a$, id quod si basis rite determinatur, semper fieri potest (art.~71). Tum index producti ex omnibus periodi terminis erit
\[= (1+2+3+\text{etc.}+t-1)a = \frac{t(t-1)}{2} a\]
i.e. $\equiv 0 \pmod{p-1}$, quando $t$ impar, et $\equiv \frac{p-1}{2}$, quando $t$ par; hinc in priori casu productum illud $\equiv +1 \pmod{p}$; in posteriori vero $\equiv -1 \pmod{p}$, (art.~62).

Q.~E.~D.

\setcounter{satz}{76}
\addtocounter{satz}{-1}
\textbf{76.}

Si numerus iste in theor. praecedente est radix primitiva, eius periodus omnes numeros $1$, $2$, $3$, \ldots, $p-1$ comprehendet, quorum productum itaque semper $\equiv -1$ (namque $p-1$ semper par, unico casu $p=2$ excepto in quo $-1 \equiv +1$ aequivalent). Theorema hoc elegans quod ita enunciari solet: productum ex omnibus numeris numero primo dato minoribus, unitate auctum per hunc primum est divisibile, primum a cel. Waring est prolatum armigeroque Wilson adscriptum, Meditt. algebr. Ed.~3.~p.~380. Sed neuter demonstrare potuit, et cel. Waring fatetur demonstrationem eo difficiliorem videri, quod nulla notatio fingi possit, quae numerum primum exprimat. At nostro quidem iudicio huiusmodi veritates ex notionibus potius quam ex notationibus hauriri debebant. Postea ill. La Grange demonstrationem dedit, Nouv. Mém. de l'Ac. de Berlin, 1771. Innititur ea considerationi coefficientium ex evolutione producti
\[(x+1)(x+2)(x+3)\ldots(x+p-1)\]
oriundorum. Scilicet posito hoc producto
\[= x^{p-1} + Ax^{p-2} + Bx^{p-3} + \ldots + Mx^2 + Nx \text{ etc.}\]
coefficientes $A$, $B$, \ldots, $M$ per $p$ erunt divisibiles, $N$ vero $=1 \cdot 2 \cdot 3 \ldots (p-1)$. Iam pro $x=-1$ productum per $p$ divisibile; tunc autem $x^{p-1}-1$ divisibile; erit $1+N \equiv 0 \pmod{p}$, quare necessario $1+N$ per $p$ dividi potest.

Denique ill. Euler in Opusc. analyt. T.~I.~p.~329 demonstrationem dedit, cum ea quam nos hic exposuimus conspirantem. Quodsi tales viri theorema hoc meditationibus suis non indignum censuerunt, non improbatum iri speramus, si aliam adhuc demonstrationem apponimus.

\setcounter{satz}{77}
\addtocounter{satz}{-1}
\textbf{77.}

Quando secundum modulum $p$, productum duorum numerorum $a$, $b$ unitati est congruum, numeros $a$, $b$ cum ill. Euler \emph{socios} vocemus. Tum secundum sect. praec. quivis numerus positivus ipso $p$ minor socium habebit positivum ipso $p$ minorem et quidem unicum. Facile autem probari potest ex numeris $1$, $2$, $3$, \ldots, $p-1$; $1$ et $p-1$ esse unicos qui sibi ipsis sint socii: numeri enim sibi ipsis socii, radices erunt congruentiae $xx \equiv 1$, quae quoniam est secundi gradus, plures quam duas radices, i.~e. alias quam $1$ et $p-1$ habere nequit. Abiectis itaque bis numerorum reliquorum $2$, $3$, \ldots, $p-2$ bini semper erunt associati; quare productum ex ipsis erit $\equiv +1$ adeoque productum ex omnibus $1$, $2$, $3$, \ldots, $p-1$, sive $p-1$, $\equiv -1$. Q.~E.~D.

Ex. gr. pro $p=13$ numeri $2$, $3$, $4$, \ldots, $11$ ita associantur: $2$ cum $7$; $3$ cum $9$; $4$ cum $10$; $5$ cum $8$; $6$ cum $11$; scilicet $2 \cdot 7 \equiv 1$; $3 \cdot 9 \equiv 1$ etc. Hinc $2 \cdot 3 \cdot 4 \ldots 11 \equiv 1$; adeoque $1 \cdot 2 \cdot 3 \ldots 12 \equiv -1$.

\setcounter{satz}{78}
\addtocounter{satz}{-1}
\textbf{78.}

Potest autem theorema Wilsonianum generalius sic proponi. Productum ex omnibus numeris, numero quocunque dato $A$ minoribus simulque ad ipsum primis, congruum est secundum $A$ unitati vel negative vel positive sumtae. Negative sumenda est unitas, quando $A$ est formae $p^\alpha$, aut huiusce $2p^\alpha$, designante $p$ numerum primum a $2$ diversum, insuperque quando $A=4$; positive autem in omnibus casibus reliquis. Theorema, quale a cel. Wilson est prolatum, sub casu priori continetur. --- Ex. gr. pro $A=15$ productum e numeris $1$, $2$, $4$, $7$, $8$, $11$, $13$, $14$ est $\equiv -1 \pmod{15}$. Demonstrationem brevitatis gratia non adiungimus: observamus tantum, eam simili modo perfici posse ut in art. praec., excepto quod congruentia $xx \equiv 1$ plures quam duas radices habere potest, quae considerationes quasdam peculiares postulant. Posset etiam demonstratio ex consideratione indicum peti, similiter ut in art.~75, si ea quae mox de modulis non primis trademus, conferantur.

\setcounter{satz}{79}
\addtocounter{satz}{-1}
\textbf{79.}

Revertimur ad enumerationem aliarum propositionum (art.~75).

Summa omnium terminorum periodi numeri cuiusvis est $\equiv 0$, uti in ex. art.~75, $1+5+12+8=26 \equiv 0 \pmod{13}$.

\textsc{Dem.} Numerus de cuius periodo agitur, sit $=a$ atque exponens ad quem pertinet, $=t$, eritque summa terminorum omnium periodi,
\[1 + a + aa + a^3 + \text{etc.} + a^{t-1} \equiv \frac{a^t-1}{a-1} \pmod{p}.\]
At $a^t-1 \equiv 0$; quare summa haec semper $\equiv 0$, nisi forte $a-1$ per $p$ sit divisibilis, sive $a \equiv 1$; hunc igitur casum excipere oportet, vel unum terminum periodum vocare velimus.

\setcounter{satz}{80}
\addtocounter{satz}{-1}
\textbf{80.}

Productum ex omnibus radicibus primitivis est $\equiv +1$, excepto unico casu, $p=3$, tum enim una tantum datur radix primitiva $2 \equiv -1$.

\textsc{Demonstr.} Si radix primitiva quaecunque pro basi assumitur, indices radicum omnium primitivarum erunt numeri ad $p-1$ primi simulque ipso minores. At horum numerorum summa, i.~e. index producti ex omnibus radicibus primitivis, est $\equiv 0 \pmod{p-1}$ adeoque productum $\equiv +1 \pmod{p}$; facile enim perspicitur, si $k$ fuerit numerus ad $p-1$ primus, etiam $p-1-k$ ad $p-1$ primum fore adeoque binos numeros ad $p-1$ primos summam constituere per $p-1$ divisibilem; ($k$ autem ipsi $p-1-k$ numquam aequalis esse potest, praeter casum, $p-1-k=2$, sive $p=3$, quem excepimus; manifesto enim $k=\frac{p-1}{2}$ in omnibus reliquis casibus $p-1$ non est primus). Q.~E.~D.

\setcounter{satz}{81}
\addtocounter{satz}{-1}
\textbf{81.}

Summa omnium radicum primitivarum est aut $\equiv 0$ (quando $p-1$ per quadratum aliquod est divisibilis), aut $\equiv +1$ vel $\equiv -1 \pmod{p}$ (quando $p-1$ est productum e numeris primis inaequalibus; quorum multitudo si est par signum positivum, si vero impar, negativum sumendum).

\textsc{Ex.} $1^\circ$ pro $p=13$, habentur radices primitivae $2$, $6$, $7$, $11$, quarum summa $26 \equiv 0 \pmod{13}$.

$2^\circ$ pro $p=11$, radices primitivae sunt $2$, $6$, $7$, $8$ quorum summa $23 \equiv +1 \pmod{11}$.

$3^\circ$ pro $p=31$, radices primitivae sunt $3$, $11$, $12$, $13$, $17$, $21$, $22$, $24$, quorum summa $123 \equiv -1 \pmod{31}$.

\textsc{Demonstr.} Supra demonstravimus (art.~55, II), si $p-1$ fuerit $=a^\alpha b^\beta c^\gamma$ etc. (designantibus $a, b, c$ etc. numeros primos inaequales) atque $A$, $B$, $C$ etc. numeri quicunque ad exponentes $a^\alpha$, $b^\beta$, $c^\gamma$ etc. respective pertinentes, omnia producta $ABC$ etc. exhibere radices primitivas. Facile vero etiam demonstrari potest, quamvis radicem primitivam per huiusmodi productum exhiberi posse et quidem unico tantum modo\footnote{Determinentur scilicet numeri $a$, $b$, $c$ etc. ita, ut sit $a \equiv 1 \pmod{a^\alpha}$ et $\equiv 0 \pmod{b^\beta c^\gamma \text{ etc.}}$, $b \equiv 1 \pmod{b^\beta}$ et $\equiv 0 \pmod{a^\alpha c^\gamma \text{ etc.}}$ etc. (vid. art.~32), unde fiet $a+b+c+\text{etc.} \equiv 1 \pmod{p-1}$ (art.~19). Iam si radix primitiva quaecunque, $r$, per productum $ABC$ etc. exhiberi debet, accipiatur $A=r^a$, $B=r^b$; $C=r^c$ etc., atque pertinebunt $A$ ad exponentem $a^\alpha$, $B$ ad exponentem $b^\beta$ etc.; productum ex omnibus $A$, $B$, $C$ etc. erit $\equiv r \pmod{p}$; denique facile perspicitur $A$, $B$, $C$ etc. alio modo determinari non posse.}. Unde sequitur haec producta loco ipsarum radicum primitivarum accipi posse. At quoniam in his productis omnes valores ipsius $A$ cum omnibus ipsius $B$ etc. combinari oportet, omnium horum productorum summa aequalis est producto ex summa omnium valorum ipsius $A$ in summam omnium valorum ipsius $B$ in summam omnium valorum ipsius $C$ etc., uti ex doctrina combinationum notum est. Designentur omnes valores ipsorum $A$, $B$ etc., per $A$, $A'$, $A''$ etc., $B$, $B'$, $B''$ etc. etc., eritque summa omnium radicum primitivarum
\[= (A+A'+A''+\text{etc.})(B+B'+\text{etc.})\text{ etc.}\]
Iam dico, si exponens $a^\alpha=1$, summam $A+A'+A''+\text{etc.}$ fore $\equiv -1 \pmod{p}$; si vero $a^\alpha > 1$, summam hanc fore $\equiv 0$, similiterque de reliquis $b^\beta$, $c^\gamma$ etc. Simulac haec erunt demonstrata, theorematis nostri veritas manifesta erit. Quando enim $p-1$ per quadratum aliquod divisibilis est, aliquis exponentium $a^\alpha$, $b^\beta$, $c^\gamma$ etc. unitatem superabit, adeoque aliquis factorum, quorum producto congrua est summa omnium radicum primitivarum erit $\equiv 0$, et proin etiam productum ipsum; quando vero $p-1$ per nullum quadratum dividi potest, omnes exponentes $a^\alpha$, $b^\beta$, $c^\gamma$ etc. erunt $=1$, unde summa omnium radicum primitivarum congrua erit producto ex tot factoribus, quorum quisque $\equiv -1$, quot habentur numeri $a$, $b$, $c$ etc., adeoque erit $\equiv +1$ vel $\equiv -1$, prout horum numerorum multitudo par vel impar. Illa autem ita probantur.

$1^\circ$. Quando $a^\alpha=1$ atque $A$ numerus ad exponentem $a^\alpha$ pertinens, reliqui numeri ad hunc exponentem pertinentes erunt $A^2$, $A^3$, \ldots, $A^{a-1}$. At $A+A^2+A^3+\ldots+A^{a-1}$ est summa periodi completae, adeoque $\equiv -1$ (art.~79), quare $A+A^2+A^3+\ldots+A^{a-1} \equiv -1$.

$2^\circ$. Quando autem $a^\alpha > 1$, atque $A$ numerus ad exponentem $a^\alpha$ pertinens, reliqui numeri ad hunc exponentem pertinentes habebuntur, si ex his $A^2$, $A^3$, \ldots, $A^{a^\alpha-1}$ reiiciuntur $A^a$, $A^{2a}$, $A^{3a}$ etc., vid. art.~53; quare summa eorum erit
\[= A+A^2+\ldots+A^{a^\alpha-1} - (A^a+A^{2a}+A^{3a}+\ldots+A^{a^\alpha-a})\]
i.~e. congrua differentiae duarum periodorum, adeoque $\equiv 0$. Q.~E.~D.


\subsection*{De modulis qui sunt numerorum primorum potestates.}

\setcounter{satz}{82}
\addtocounter{satz}{-1}
\textbf{82.}

Omnia quae hactenus exposuimus innituntur suppositioni modulum esse numerum primum. Superest ut eum quoque casum consideremus, ubi pro modulo assumitur numerus compositus. Attamen quum hic neque proprietates tam elegantes eniteant, quam in casu priori, neque ad eas inveniendas artificiis subtilibus sit opus, sed potius omnia fere per solam principiorum praecedentium applicationem erui possint, omnes minutias hic exhaurire superfluum atque taediosum foret. Breviter itaque quae huic casui cum priori sint communia quaeque propria exponemus.

\setcounter{satz}{83}
\addtocounter{satz}{-1}
\textbf{83.}

Propositiones artt.~45---48 generaliter iam sunt demonstratae. At prop. art.~49 ita immutari debet:

Si $f$ designat, quot numeri dentur ad $m$ primi simul ipso $m$ minores, i.~e. si $f = \varphi m$ (art.~38): exponens $t$ infimae potestatis numeri dati $a$ ad $m$ primi, quae secundum modulum $m$ unitati est congrua, vel erit $=f$ vel pars aliquota huius numeri.

Demonstratio prop. art.~49 etiam pro hoc casu valere potest, si modo ubique loco ipsius $p$, $m$, loco ipsius $p-1$, $f$, et loco numerorum $1$, $2$, $3$, \ldots, $p-1$, numeri ad $m$ primi simulque ipso $m$ minores substituantur. Huc itaque lectorem ablegamus. Ceterum demonstrationes reliquae de quibus illic locuti sumus (artt.~50,~51), non sine multis ambagibus ad hunc casum applicari possunt. --- At respectu propositionum sequentium, art.~52 sqq. magna differentia incipit inter modulos, qui numerorum primorum sunt potestates, eosque, qui per plures numeros primos dividi possunt. Seorsim itaque modulos prioris generis contemplabimur.

\setcounter{satz}{84}
\addtocounter{satz}{-1}
\textbf{84.}

Si modulus $m=p^n$, designante $p$ numerum primum, erit $f=p^{n-1}(p-1)$ (art.~38). Iam si disquisitiones in artt.~53,~54 contentae ad hunc casum applicantur, mutatis mutandis uti in art. praec. praescripsimus, invenietur, omnia quae ibi demonstrata sunt etiam pro hoc casu locum habere, si modo ante probatum esset, congruentiam formae $x^t \equiv 1 \pmod{p^n}$ plures quam $t$ radices diversas habere non posse. Pro modulo primo hanc veritatem ex propositione generaliori art.~43 deduximus, quae autem in omni sua extensione de modulis primis tantummodo valet, neque adeo ad hunc casum applicanda. Attamen propositionem pro hoc casu particulari veram esse, per methodum singularem demonstrabimus. Infra (sect.~VIII) idem facilius invenire docebimus.

\setcounter{satz}{85}
\addtocounter{satz}{-1}
\textbf{85.}

Demonstrandum proponimus nobis hoc theorema:

Si numerorum $t$ et $p^{n-1}(p-1)$ divisor communis maximus est $e$, congruentia $x^t \equiv 1 \pmod{p^n}$ habebit $e$ radices diversas.

Sit $e=kp^\mu$, ita ut $k$ factorem $p$ non involvat, adeoque numerum $p-1$ metiatur. Tum congruentia $x^t \equiv 1$ secundum modulum $p$ habebit $k$ radices diversas, quibus per $A$, $B$, $C$ etc. designatis radix quaecunque eiusdem congruentiae secundum modulum $p^n$, congrua esse debet secundum modulum $p$ alicui numerorum $A$, $B$, $C$ etc. Iam demonstrabimus, congruentiam $x^t \equiv 1 \pmod{p^n}$ habere $p^\mu$ radices ipsi $A$, totidem ipsi $B$ etc. congruas secundum modulum $p$. Quo facto omnium radicum numerus erit $kp^\mu$ sive $e$, uti diximus. Illam vero demonstrationem ita adornabimus, ut primo ostendamus, si $a$ fuerit radix ipsi $A$ secundum modulum $p$ congrua, etiam $a+p^{n-1}$, $a+2p^{n-1}$, $a+3p^{n-1}$, \ldots, $a+(p-1)p^{n-1}$ fore radices; secundo, numeros ipsi $A$ secundum modulum $p$ congruos alios quam qui in forma $a+hp^{n-1}$ sint comprehensi (denotante $h$ integrum quemcunque), radices esse non posse: unde manifesto $p^\mu$ radices diversae habebuntur, et non plures: atque idem etiam de radicibus, quae singulis $B$, $C$ etc. sunt congruae, locum habebit: tertio docebimus, quomodo semper radix, ipsi $A$ secundum $p$ congrua, inveniri possit.

\setcounter{satz}{86}
\addtocounter{satz}{-1}
\textbf{86.}

\textsc{Theorema.} Si uti in art. praec. $t$ est numerus per $p^\mu$, neque vero per $p^{\mu+1}$ divisibilis, erit
\[(a+hp^\nu)^t \equiv a^t \pmod{p^{\nu+\mu+1}}\quad\text{at}\quad (a+hp^\nu)^t \not\equiv a^t \pmod{p^{\nu+\mu}}.\]

Theorematis pars posterior locum non habet, quando $p=2$ simulque $\nu=1$.

Demonstratio huius theorematis ex evolutione potestatis binomii peti posset, si ostenderetur omnes terminos post secundum per $p^{\nu+\mu+1}$ divisibiles esse. Sed quoniam consideratio denominatorum coefficientium in aliquot ambages deducit, methodum sequentem praeferimus.

Ponamus primo $[a+b]^t \equiv 1$ atque $v \equiv 1$, eritque propter
\[x^t-y^t = (x-y)(x^{t-1}+x^{t-2}y+x^{t-3}y^2+\text{etc.}+y^{t-1})\]
\[(a+hp^\nu)^t - a^t = hp^\nu\bigl((a+hp^\nu)^{t-1} + (a+hp^\nu)^{t-2}a + \text{etc.} + a^{t-1}\bigr).\]
At est $a+hp^\nu \equiv a \pmod{p^\nu}$, quare quisque terminus $(a+hp^\nu)^{t-1}$, $(a+hp^\nu)^{t-2}a$ etc. erit $\equiv a^{t-1} \pmod{p^\nu}$, adeoque omnium summa $\equiv ta^{t-1} \pmod{p^\nu}$ sive formae $ta^{t-1}+Vp^\nu$ denotante $V$ numerum quemcunque. Hinc $(a+hp^\nu)^t - a^t$ erit formae $a^{t-1}hp^\nu t + Vhp^{2\nu}$, i.e. $\equiv ta^{t-1}hp^\nu \pmod{p^{\nu+\mu}}$ et $\equiv 0 \pmod{p^{\nu+\mu+1}}$. Pro hoc itaque casu theorema est demonstratum.

Iam si theorema pro aliis ipsius $v$ valoribus verum non esset, manente etiamnum $\mu \geq 1$, limes aliquis necessario daretur, usque ad quem theorema semper verum foret, ultra vero falsum. Sit minimus valor ipsius $v$, pro quo falsum est $=q$, unde facile perspicitur, si $t$ per $p^{q-1}$ non autem per $p^q$ fuerit divisibilis, theorema adhuc verum esse, at si loco ipsius $t$ substituatur $tp$, falsum. Habemus itaque
\[(a+hp^\nu)^t \equiv a^t + a^{t-1}hp^\nu t \pmod{p^{\nu+q}}\quad\text{sive}\quad = a^t + a^{t-1}hp^\nu t + up^{\nu+q}\]
denotante $u$ numerum integrum. At quia pro $v=1$ theorema iam est demonstratum, erit
\[(a^t + a^{t-1}hp^\nu t + up^{\nu+q})^p \equiv a^{tp} + a^{tp-1}hp^{\nu+1}tp + a^{tp-1}up^{\nu+q+1} \pmod{p^{\nu+q+1}}\]
adeoque etiam
\[(a+hp^\nu)^{tp} \equiv a^{tp} + a^{tp-1}hp^{\nu+1}tp \pmod{p^{\nu+q+1}}\]
i.~e. theorema etiam verum, si loco ipsius $t$ substituitur $tp$, i.~e. etiam pro $v=q$, contra hypothesin. Unde manifestum pro omnibus ipsius $v$ valoribus theorema verum esse.

\setcounter{satz}{87}
\addtocounter{satz}{-1}
\textbf{87.}

Superest casus ubi $\mu=1$. Per methodum prorsus similem ei qua in art. praec. usi sumus, sine adiumento theorematis binomialis demonstrari potest, esse
\begin{align*}
(a+hp)^t - a^t &\equiv a^{t-1}hp \cdot t \pmod{p^{\nu+1}} \\
a(a+hp)^{t-1} &\equiv a^t + a^{t-1}hp(t-1) \pmod{p^{\nu+1}} \\
a^2(a+hp)^{t-2} &\equiv a^t + a^{t-1}hp(t-2) \pmod{p^{\nu+1}} \\
&\text{etc.}
\end{align*}
unde aggregatum erit (quia partium multitudo $=t$) $\equiv ta^t + a^{t-1}hp \cdot \frac{t(t-1)}{2} \pmod{p^{\nu+1}}$.

At quoniam $t$ per $p$ divisibilis, etiam $\frac{t(t-1)}{2}$ per $p$ divisibilis erit in omnibus casibus excepto eo ubi $p=2$ de quo iam in art. praec. monuimus. In reliquis autem casibus erit $a^{t-1}hp \cdot \frac{t(t-1)}{2} \equiv 0 \pmod{p^{\nu+1}}$, adeoque etiam illud aggregatum $\equiv ta^t \pmod{p^{\nu+1}}$ ut in art. praec. In reliquis demonstratio hic eodem modo procedit ut istic.

Colligimus igitur generaliter, unico casu $p=2$ excepto, esse $(a+hp^\nu)^t \equiv a^t \pmod{p^{\nu+\mu}}$ et $(a+hp^\nu)^t \not\equiv a^t$ pro quovis modulo qui sit altior potestas ipsius $p$, quam haec $p^{\nu+\mu}$, quoties quidem $h$ per $p$ non est divisibilis, atque $p^\mu$ potestas suprema ipsius $p$ quae numerum $t$ dividit.

Hinc protinus derivantur propositiones 1. et 2, quas art.~85 demonstrandas nobis proposueramus: scilicet primo, si $a^t \equiv 1$, erit etiam $(a+hp^{n-1})^t \equiv 1 \pmod{p^n}$; secundo si numerus aliquis $d$ ipsi $A$ adeoque etiam ipsi $a$ secundum modulum $p$ congruus, neque vero huic secundum modulum $p^{n-1}$, congruentiae $x^t \equiv 1 \pmod{p^n}$ satisfaceret: ponamus $d$ esse $= a+lp^{n-1}$, ita ut $l$ per $p$ non sit divisibilis, eritque $d^t \not\equiv 1$, tunc autem $(a+lp^{n-1})^t$ secundum modulum $p^{n+\mu}$ ipsi $a^t$ congruus erit, non autem secundum modulum $p^n$, quae est altior potestas, quare $d$ radix congruentiae $x^t \equiv 1$ esse nequit.

\setcounter{satz}{88}
\addtocounter{satz}{-1}
\textbf{88.}

Tertium vero fuit radicem aliquam congruentiae $x^t \equiv 1 \pmod{p^n}$, ipsi $A$ congruam, invenire. Ostendemus hic tantummodo, quomodo hoc fieri possit, si iam radix eiusdem congruentiae secundum modulum $p^{n-1}$ innotuerit; manifesto hoc sufficit, quum a modulo $p$, pro quo $A$ est radix, ad modulum $p^2$, sicque deinceps ad omnes potestates consecutivas progredi possimus.

Esto itaque $a$ radix congruentiae $x^t \equiv 1 \pmod{p^{n-1}}$, quaeriturque radix eiusdem congruentiae secundum modulum $p^n$. Ponatur haec $= a+hp^{n-\nu-1}$, quam formam eam habere debere ex art. praec. sequitur (casum ubi $v=n-1$ postea seorsim considerabimus: maior vero quam $n-1$, $v$ esse nequit). Debet itaque esse $(a+hp^{n-\nu-1})^t \equiv 1 \pmod{p^n}$.

At $(a+hp^{n-\nu-1})^t \equiv a^t + a^{t-1}htp^{n-\nu-1} \pmod{p^n}$. Si itaque $h$ ita determinatur, ut fiat $a^t + a^{t-1}htp^{n-\nu-1} \equiv 1 \pmod{p^n}$, sive (quia per hyp. $1 \equiv a^t \pmod{p^{n-1}}$ atque $t$ per $p^\nu$ divisibilis) ita ut fiat $a^{t-1}h \cdot \frac{t}{p^\nu}$ per $p$ divisibilis, quaesito satisfactum erit. Hoc autem semper fieri posse ex Sect. praec. manifestum, quum $t$ per altiorem ipsius $p$ potestatem quam $p^\nu$ dividi non posse hic supponamus, adeoque $a^{t-1}$ ad $p$ sit primus.

Si vero $v=n-1$, i.~e. $t$ per $p^{n-1}$ sive etiam per altiorem ipsius $p$ potestatem divisibilis, quivis valor $A$ congruentiae $x^t \equiv 1$ secundum modulum $p$ satisfaciens eidem etiam secundum modulum $p^n$ satisfaciet. Sit enim $t=p^{n-1}x$, eritque $t \equiv 0 \pmod{p^{n-1}}$; quare quoniam $A^t \equiv 1 \pmod{p}$, erit etiam $A^t \equiv 1 \pmod{p^n}$. Ponatur itaque $A^t = 1+hp$, eritque $A^{tp} = (1+hp)^p \equiv 1 \pmod{p^n}$ art.~87.

\setcounter{satz}{89}
\addtocounter{satz}{-1}
\textbf{89.}

Omnia quae art.~57 sqq. adiumento theorematis, congruentiam $x^n \equiv 1$ plures quam $t$ radices diversas non habere eruimus, etiam pro modulo qui est numeri primi potestas locum habent, et si radices primitivae vocantur numeri, qui ad exponentem $p^{n-1}(p-1)$ pertinent, sive in quorum periodis omnes numeri per $p$ non divisibiles inveniuntur, etiam hic radices primitivae exstabunt. Omnia autem quae supra de indicibus eorumque usu tradidimus necnon de solutione congruentiae $x^n \equiv 1$, ad hunc quoque casum applicari possunt. Quae quum nulli difficultati obnoxia sint, omnia ex integro repetere superfluum foret. Praeterea radices congruentiae $x^t \equiv 1$ secundum modulum $p^n$ e radicibus eiusdem congruentiae secundum $p$ deducere docuimus. Sed de eo casu ubi potestas aliqua numeri $2$ est modulus, quia supra exceptus fuit, aliqua adhuc sunt adiicienda.

\subsection*{Moduli qui sunt potestates binarii.}

\setcounter{satz}{90}
\addtocounter{satz}{-1}
\textbf{90.}

Si potestas aliqua numeri $2$, altior quam secunda, puta $2^n$ pro modulo accipitur, numeri cuiusvis imparis potestas exponentis $2^{n-2}$, unitati est congrua.

Ex. gr. $3^8 = 6561 \equiv 1 \pmod{32}$.

Quivis enim numerus impar vel sub forma $1+4h$, vel sub hac $-1+4h$ comprehenditur: unde propositio protinus sequitur (theor. art.~86).

Quoniam igitur exponens, ad quem quicunque numerus impar secundum modulum $2^n$ pertinet, divisor ipsius $2^{n-2}$ esse debet, quivis ad aliquem horum numerum pertinebit $1$, $2$, $4$, $8$, \ldots, $2^{n-2}$, ad quemnam vero pertineat, ita facile diiudicatur. Sit numerus propositus $=4h+1$, atque exponens maximae potestatis numeri $2$, quae ipsum $h$ metitur, $=m$ (qui etiam $=0$ esse potest, quando scilicet $h$ est impar); tum exponens, ad quem numerus propositus pertinet, erit $=2^{n-m-2}$, siquidem $n \geq m+2$; si autem $n<m+1$, vel $n=m+1$, numerus propositus est $=+1$ adeoque vel ad exponentem $1$ vel ad exponentem $2$ pertinebit. Numerum enim formae $1+2^{m+1}k$, (quae huic aequivalet, $4h+1$) ad potestatem exponentis $2^{n-m-2}$ elevatum unitati secundum modulum $2^n$ congruum fieri, ad potestatem autem exponentis, qui est inferior numeri $2$ potestas, incongruum, ex art.~86 nullo negotio deducitur. Numerus itaque quicunque formae $8h+3$ vel $8h+5$ ad exponentem $2^{n-2}$ pertinebit.

\setcounter{satz}{91}
\addtocounter{satz}{-1}
\textbf{91.}

Hinc patet eo sensu, quo supra expressionem accepimus, radices primitivas hic non dari, nullos scilicet numeros, quorum periodus omnes numeros modulo minores ad ipsumque primos amplectatur. Attamen facile perspicitur, analogon hic haberi. Invenitur enim, numeri formae $8h+3$ potestatem exponentis imparis semper esse formae $8h+3$, potestatem autem exponentis paris, semper formae $8h+1$; nulla igitur potestas formae $8h+5$ aut $8h+7$ esse potest. Quare quum periodus numeri formae $8h+5$ ex $2^{n-2}$ terminis diversis constet, quorum quisque aut formae $8h+3$ aut huius $8h+1$, neque plures huiusmodi numeri modulo minores dentur quam $2^{n-2}$, manifeste quivis numerus formae $8h+1$ vel $8h+3$ congruus est secundum modulum $2^n$ potestati alicui numeri cuiuscunque formae $8h+5$. Simili modo ostendi potest periodum numeri formae $8h+5$ comprehendere omnes numeros formarum $8h+1$ et $8h+5$. Si igitur numerus formae $8h+5$ pro basi assumitur, omnes numeri formae $8h+1$ et $8h+5$, positive, omnesque formae $8h+3$ et $8h+7$, negative sumti, indices reales nanciscentur, et quidem hic indices secundum $2^{n-2}$ congrui pro aequivalentibus sunt habendi. Hoc modo tabula nostra I intelligenda, ubi pro modulis $16$, $32$ et $64$ (namque pro modulo $8$ nulla tabula necessaria erit) semper numerum $5$ pro basi accepimus. Ex. gr. numero $19$, qui est formae $8n+3$ adeoque negative sumendus, respondet pro modulo $64$ index $7$, id quod significat esse $5^7 \equiv -19 \pmod{64}$. Numeris autem formarum $8n+1$, $8n+5$ negative, atque numeris formarum $8n+3$, $8n+7$ positive acceptis, indices quasi imaginarii tribuendi forent. Quos introducendo calculus indicum ad algorithmum perquam simplicem reduci potest. Sed quoniam, si haec ad omnem rigorem exponere vellemus, nimis longe evagari oporteret, hoc negotium ad aliam occasionem nobis reservamus, quando forsan fusius quantitatum imaginariarum theoriam, quae nostro quidem iudicio a nemine hactenus ad notiones claras est reducta, pertractare suscipiemus. Periti hunc algorithmum facile ipsi eruent: qui minus sunt exercitati, perinde tamen tabula hac uti poterunt ut ii qui recentiorum commenta de logarithmis imaginariis ignorant, logarithmis utuntur, si quidem principia supra stabilita probe tenuerint.

\subsection*{Moduli e pluribus primis compositi.}

\setcounter{satz}{92}
\addtocounter{satz}{-1}
\textbf{92.}

Secundum modulum e pluribus primis compositum tantum non omnia quae ad residua potestatum pertinent, ex theoria congruentiarum generali deduci possunt; quia vero infra congruentias quascunque secundum modulum e pluribus primis compositum ad congruentias, quarum modulus est primus aut primi potestas, reducere fusius docebimus, non est quod huic rei multum hic immoremur. Observamus tantum, bellissimam proprietatem, quae pro reliquis modulis locum habeat, quod scilicet semper exstant numeri quorum periodus omnes numeros ad modulum primos complectatur, hic deficere, excepto unico casu, quando scilicet modulus est duplum numeri primi, aut potestatis numeri primi. Si enim modulus $m$ redigitur ad formam $A^\alpha B^\beta C^\gamma$ etc., designantibus $A$, $B$, $C$ etc. numeros primos diversos, praeterea $A^{\alpha-1}(A-1)$ designatur per $a$, $B^{\beta-1}(B-1)$ per $b$ etc. denique $z$ est numerus ad $m$ primus; erit $z^a \equiv 1 \pmod{A^\alpha}$, $z^b \equiv 1 \pmod{B^\beta}$ etc. Quodsi igitur $f$ est minimus numerorum $a$, $b$, $c$ etc. dividuus communis, erit $z^f \equiv 1$ secundum omnes modulos $A^\alpha$, $B^\beta$ etc. adeoque etiam secundum $m$, cui illorum productum est aequale. At excepto casu, ubi $m$ est duplum numeri primi aut potestatis numeri primi numerorum $a$, $b$, $c$ etc. dividuus communis minimus, ipsorum producto est minor (quoniam numeri $a$, $b$, $c$ etc. inter se primi esse nequeunt sed certe divisorem $2$ communem habent). Nullius itaque numeri periodus tot terminos comprehendere potest, quot dantur numeri ad modulum primi ipsoque minores, quia horum numerus producto ex $a$, $b$, $c$ etc. est aequalis.

Ita ex. gr. pro $m=1001$ cuiusvis numeri ad $m$ primi potestas exponentis $60$ unitati est congrua, quia $60$ est dividuus communis numerorum $6$, $10$, $12$. Casus autem ubi modulus est duplum numeri primi aut duplum potestatis numeri primi, illi ubi est primus aut primi potestas, prorsus est similis.

\setcounter{satz}{93}
\addtocounter{satz}{-1}
\textbf{93.}

Scriptorum in quibus alii geometrae de argumento in hac sectione pertractato egerunt, iam passim mentio est facta. Eos tamen qui quaedam fusius, quam nobis brevitas permisit, explicata desiderant, ablegamus imprimis ad sequentes ill. Euleri commentationes, ob perspicuitatem, qua vir summus prae omnibus semper excelluit, maxime commendabiles:

\textit{Theoremata circa residua ex divisione potestatum relicta.} Comm. nov. Petr. T.~VII p.~49 sqq.

\textit{Demonstrationes circa residua ex divisione potestatum per numeros primos resultantia.} Ibid. T.~XVIII; p.~85 sqq.

Adiungi his possunt Opusculorum analyt. T.~I, dissertt.~5 et 8.


\section*{SECTIO QUARTA}
\addcontentsline{toc}{section}{SECTIO QUARTA}
\section*{DE CONGRUENTIIS SECUNDI GRADUS.}

\subsection*{Residua et non-residua quadratica.}

\setcounter{satz}{94}
\addtocounter{satz}{-1}
\textbf{94.}

\textsc{Theorema.} Numero quocunque $m$ pro modulo accepto, ex numeris $0$, $1$, $2$, $3$, \ldots, $m-1$, plures quam $\frac{1}{2}m+1$, quando $m$ est par, sive plures quam $\frac{1}{2}(m+1)$, quando $m$ est impar, quadrato congrui fieri non possunt.

\textsc{Dem.} Quoniam numerorum congruorum quadrata sunt congrua: quivis numerus, qui ulli quadrato congruus fieri potest, etiam quadrato alicui cuius radix $<\frac{1}{2}m$ congruus erit. Sufficit itaque residua minima quadratorum $0$, $1$, $4$, $9$, \ldots, $(m-1)^2$ considerare. At facile perspicitur, esse $(m-1)^2 \equiv 1^2$, $(m-2)^2 \equiv 2^2$, $(m-3)^2 \equiv 3^2$ etc. Hinc etiam, quando $m$ est par, quadratorum $(\frac{1}{2}m-1)^2$ et $(\frac{1}{2}m+1)^2$, $(\frac{1}{2}m-2)^2$ et $(\frac{1}{2}m+2)^2$ etc. residua minima eadem erunt: quando vero $m$ est impar, quadrata $(\frac{m-1}{2}-1)^2$ et $(\frac{m+1}{2}+1)^2$, $(\frac{m-1}{2}-2)^2$ et $(\frac{m+1}{2}+2)^2$ etc. erunt congrua. Unde palam est, alios numeros, quam qui alicui ex quadratis $0$, $1$, $4$, $9$, \ldots, $(\frac{1}{2}m)^2$ congrui sint, quadrato congruos fieri non posse, quando $m$ par; quando vero impar, quemvis numerum, qui ulli quadrato sit congruus, alicui ex his $0$, $1$, $4$, $9$, \ldots, $(\frac{m-1}{2})^2$ necessario congruum esse. Quare dabuntur ad summum in priori casu $\frac{1}{2}m+1$ residua minima diversa, in posteriori $\frac{1}{2}(m+1)$. Q.~E.~D.

\textsc{Exemplum.} Secundum modulum $13$ quadratorum numerorum $0$, $1$, $2$, $3$, \ldots, $6$ residua minima inveniuntur $0$, $1$, $4$, $9$, $3$, $12$, $10$, post haec vero eadem ordine inverso recurrunt $10$, $12$, $3$ etc. Quare numerus quisque, nulli ex istis residuis congruus, sive qui alicui ex his est congruus, $2$, $5$, $6$, $7$, $8$, $11$, nulli quadrato congruus esse potest.

Secundum modulum $15$ haec inveniuntur residua $0$, $1$, $4$, $9$, $1$, $10$, $6$, $4$, post quae eadem ordine inverso recurrunt. Hic igitur numerus residuorum, quae quadrato congrua fieri possunt, minor adhuc est quam $\frac{1}{2}m+\frac{1}{2}$, quum sint $0$, $1$, $4$, $6$, $9$, $10$. Numeri autem $2$, $3$, $5$, $7$, $8$, $11$, $12$, $13$, $14$ et qui horum alicui sunt congrui, nulli quadrato secundum mod.~$15$ congrui fieri possunt.

\setcounter{satz}{95}
\addtocounter{satz}{-1}
\textbf{95.}

Hinc colligitur, pro quovis modulo omnes numeros in duas classes distingui posse, quarum altera contineat numeros, qui quadrato alicui congrui fieri possint, altera eos, qui non possint. Illos appellabimus \emph{residua quadratica} numeri istius quem pro modulo accepimus\footnote{Proprie quidem hic casu secundo alio sensu utimur, quam hucusque fecimus. Dicere scilicet oporteret, $r$ esse residuum quadrati $aa$ secundum modulum $m$ quando $r \equiv aa \pmod{m}$; at brevitatis gratia in hac sectione semper $r$ ipsius $m$ residuum quadraticum vocamus neque hinc ulla ambiguitas metuenda. Expressionem enim, residuum, quando idem significat quod numerus congruus, abhinc non adhibebimus, nisi forte de residuis minimis sermo sit, ubi nullum dubium esse potest.}, hos vero ipsius \emph{non-residua quadratica}, sive etiam, quoties ambiguitas nulla inde oriri potest, simpliciter \emph{residua} et \emph{non-residua}. Ceterum palam est sufficere, si omnes numeri $0$, $1$, $2$, \ldots, $m-1$ in classes redacti sint: numeri enim congrui ad eandem classem erunt referendi.

Etiam in hac disquisitione a modulis primis initium faciemus, quod itaque subintelligendum erit, etiam si expressis verbis non moneatur. Numerus primus $2$ autem excludendus, sive numeri primi impares tantum considerandi.

\subsection*{Quoties modulus est numerus primus, multitudo residuorum ipso minorum multitudini non-residuorum aequalis.}

\setcounter{satz}{96}
\addtocounter{satz}{-1}
\textbf{96.}

Numero primo $p$ pro modulo accepto, numerorum $1$, $2$, $3$, \ldots, $p-1$ semissis erunt residua quadratica, reliqui non-residua, i.~e. dabuntur $\frac{1}{2}(p-1)$ residua totidemque non-residua.

Facile enim probatur, omnia quadrata $1$, $4$, $9$, \ldots, $(\frac{p-1}{2})^2$ esse incongrua. Scilicet si fieri posset $r^2 \equiv r'^2 \pmod{p}$ atque numeri $r$, $r'$ inaequales et non maiores quam $\frac{1}{2}(p-1)$ posito i.~q. licet $r>r'$, fieret $(r-r')(r+r')$ positivus et per $p$ divisibilis. At uterque factor $r-r'$, et $r+r'$ ipso $p$ est minor, quare suppositio consistere nequit (art.~13). Habentur itaque $\frac{1}{2}(p-1)$ residua quadratica inter hos numeros $1$, $2$, $3$, \ldots, $p-1$ contenta; plura vero inter ipsos esse nequeunt, quia accedente residuo $0$ prodeunt $\frac{1}{2}(p+1)$, quem numerum omnium residuorum multitudo superare nequit. Quare reliqui numeri erunt non-residua horumque multitudo $=\frac{1}{2}(p-1)$.

Quum cifra semper sit residuum, hanc numerosque per modulum divisibiles ab investigationibus his excludimus, quia hic casus per se est clarus, theorematumque concinnitatem tantum turbaret. Ex eadem caussa etiam modulum $2$ exclusimus.

\setcounter{satz}{97}
\addtocounter{satz}{-1}
\textbf{97.}

Quum plura quae in hac Sect. exponemus, etiam ex principiis Sect. praec. derivari possint, neque inutile sit, eandem veritatem per methodos diversas perscrutari, hunc nexum ostendemus. Facile vero intelligitur, omnes numeros quadrato congruos, indices pares habere, eos contra, qui quadrato nullo modo congrui fieri possint, impares. Quia vero $p-1$ est numerus par, tot indices pares erunt quot impares, scilicet $\frac{1}{2}(p-1)$, totidemque tum residua tum non-residua dabuntur.

\textsc{Exempla.} Pro modulis

\begin{tabular}{c|l}
$3$ & sunt residua $1$ \\
$5$ & $1$, $4$ \\
$7$ & $1$, $2$, $4$ \\
$11$ & $1$, $3$, $4$, $5$, $9$ \\
$13$ & $1$, $3$, $4$, $9$, $10$, $12$ \\
$17$ & $1$, $2$, $4$, $8$, $9$, $13$, $15$, $16$ etc.
\end{tabular}

reliqui vero numeri his modulis minores, non-residua.

\subsection*{Quaestio, utrum numerus compositus residuum numeri primi dati sit an non-residuum, ab indole factorum pendet.}

\setcounter{satz}{98}
\addtocounter{satz}{-1}
\textbf{98.}

\textsc{Theorema.} Productum e duobus residuis quadraticis numeri primi $p$, est residuum; productum e residuo in non-residuum, est non-residuum; denique productum e duobus non-residuis, residuum.

\textsc{Demonstr.} I. Sint $A$, $B$ residua e quadratis $aa$, $bb$ oriunda sive $A \equiv aa$, $B \equiv bb$, eritque productum $AB$ quadrato numeri $ab$ congruum e. residuum.

II. Quando $A$ est residuum puta $A \equiv aa$, $B$ vero non-residuum, $AB$ erit non-residuum. Ponatur enim, si fieri potest, $AB \equiv kk$, sitque valor expressionis $\frac{k}{a} \pmod{p} = b$; erit itaque $aaB \equiv aabb$, unde $B \equiv bb$, i.~e. $B$ residuum contra hyp.

Aliter. Multiplicentur omnes numeri, qui inter hos $1$, $2$, $3$, \ldots, $p-1$ sunt residua (quorum multitudo $=\frac{1}{2}(p-1)$), per $A$, omniaque producta erunt residua quadratica, et quidem erunt omnia incongrua. Iam si non-residuum $B$ per $A$ multiplicatur, productum nulli productorum quae iam habentur congruum erit; quare si residuum esset, haberentur $\frac{1}{2}(p+1)$ residua incongrua, inter quae nondum est residuum $0$, contra art.~96.

III. Sint $A$, $B$ non-residua. Multiplicentur omnes numeri, qui inter hos $1$, $2$, $3$, \ldots, $p-1$ sunt residua, per $A$, habebunturque $\frac{1}{2}(p-1)$ non-residua inter se incongrua (II); iam productum $AB$ nulli illorum congruum esse potest; quodsi igitur esset non-residuum, haberentur $\frac{1}{2}(p+1)$ non-residua inter se incongrua, contra art.~96. Quare productum etc. Q.~E.~D.

Facilius adhuc haec theoremata e principiis sect. praec. derivantur. Quia enim residuorum indices semper sunt pares, non-residuorum vero impares, index producti e duobus residuis vel non-residuis erit par, adeoque productum ipsum, residuum. Contra index producti e residuo in non-residuum erit impar adeoque productum ipsum non-residuum.

Utraque demonstrandi methodus etiam pro his theorematibus adhiberi potest: Expressionis $\frac{a}{b} \pmod{p}$ valor erit residuum, quando numeri $a$, $b$ simul sunt residua, vel simul non-residua; contra autem erit non-residuum, quando numerorum $a$, $b$ alter est residuum, alter non-residuum. Possunt etiam ex conversione theorr. praec. obtineri.

\setcounter{satz}{99}
\addtocounter{satz}{-1}
\textbf{99.}

Generaliter, productum ex quotcunque factoribus est residuum tum quando omnes sunt residua, tum quando non-residuorum, quae inter eos occurrunt, multitudo est par; quando vero multitudo non-residuorum quae inter factores reperiuntur, est impar, productum erit non-residuum. Facile itaque diiudicari potest, utrum numerus compositus sit residuum necne, si modo, quid sint singuli ipsius factores constet. Quamobrem in tabula II numeros primos tantummodo recepimus.

Oeconomia huius tabulae haec est. In margine positi sunt moduli\footnote{Quomodo etiam modulis compositis carere possimus, mox docebimus.}, in facie vero numeri primi successivi; quando ex his aliquis fuit residuum moduli alicuius, in spatio utrique respondente lineola collocata est, quando vero numerus primus fuit non-residuum moduli, spatium respondens vacuum mansit.

\subsection*{De modulis, qui sunt numeri compositi.}

\setcounter{satz}{100}
\addtocounter{satz}{-1}
\textbf{100.}

Antequam ad difficiliora progrediamur, quaedam de modulis non primis adiicienda sunt.

Si numeri primi $p$, potestas aliqua $p^n$ pro modulo assumitur (ubi $p$ non esse $2$ supponimus), omnium numerorum per $p$ non divisibilium moduloque minorum altera semissis erunt residua, altera non-residua, i.~e. utrorumque multitudo $=\frac{1}{2}(p-1)p^{n-1}$.

Si enim $r$ est residuum: quadrato alicui congruus erit, cuius radix moduli dimidium non superat, vid. art.~94. Iam facile perspicitur, dari $\frac{1}{2}(p-1)p^{n-1}$ numeros per $p$ non divisibiles modulique semisse minores; superest itaque ut demonstretur omnium horum numerorum quadrata incongrua esse, sive residua quadratica diversa suppeditare. Quodsi duorum numerorum $a$, $b$ per $p$ non divisibilium modulique semisse minorum quadrata essent congrua, foret $aa \equiv bb$ sive $(a-b)(a+b)$ per $p^n$ divisibilis (posito i.~q. licet $a>b$). Hoc vero fieri non potest, nisi vel alter numerorum $a-b$, $a+b$ per $p^n$ fuerit divisibilis, quod fieri nequit, quoniam uterque $<p^n$, vel alter per $p$, alter vero per $p^{n-1}$, i.~e. uterque per $p$. Sed etiam hoc fieri nequit. Manifesto enim etiam summa et differentia $2a$ et $2b$ per $p$ foret divisibilis adeoque etiam $a$ et $b$ contra hyp. Hinc tandem colligitur inter numeros per $p$ non divisibiles moduloque minores $\frac{1}{2}(p-1)p^{n-1}$ residua dari, reliquos quorum multitudo aeque magna, esse non-residua. Q.~E.~D. --- Potest etiam theorema hoc ex consideratione indicum derivari simili modo ut art.~97.

\setcounter{satz}{101}
\addtocounter{satz}{-1}
\textbf{101.}

Quivis numerus per $p$ non divisibilis, qui ipsius $p$ est residuum, erit residuum etiam ipsius $p^n$; qui vero ipsius $p$ est non-residuum, etiam ipsius $p^n$ non-residuum erit.

Pars posterior huius propositionis per se est manifesta. Si itaque prior falsa esset, inter numeros ipso $p^n$ minores simulque per $p$ non divisibiles plures forent residua ipsius $p$ quam ipsius $p^n$, i.~quam $\frac{1}{2}p^{n-1}(p-1)$. Nullo vero negotio perspici poterit, multitudinem residuorum numeri $p$ inter illos numeros esse praecise $=\frac{1}{2}p^{n-1}(p-1)$.

Aeque facile est, quadratum reipsa invenire, quod secundum modulum $p^n$ residuo dato sit congruum, si quadratum huic residuo secundum modulum $p$ congruum habetur.

Scilicet si quadratum habetur, $aa$, quod residuo dato $A$ secundum modulum $p^\alpha$ est congruum, deducitur inde quadratum ipsi $A$ secundum modulum $p^{2\alpha}$ congruum (ubi $\nu \geq \alpha$ et $=\nu$ vel $<2\alpha$ supponitur) sequenti modo. Ponatur radix quadrati quaesiti $= \pm a + xp^\nu$, quam formam eam habere debere facile perspicitur; debetque esse $aa \pm 2axp^\nu + xx p^{2\nu} \equiv A \pmod{p^{2\alpha}}$ sive propter $2\nu > \nu$, $A - aa \equiv \pm 2axp^\nu \pmod{p^{2\alpha}}$. Sit $A - aa = p^\nu d$, eritque valor expressionis $\frac{d}{\pm 2a} \pmod{p^{2\alpha-\nu}}$, quae huic $\frac{d}{\pm 2a} \pmod{p^{\nu}}$ aequivalet.

Dato igitur quadrato ipsi $A$ secundum $p$ congruo deducitur inde quadratum ipsi $A$ secundum modulum $p^2$ congruum hinc ad modulum $p^4$, hinc ad $p^8$ etc. ascendi poterit.

\textsc{Ex.} Proposito residuo $6$, quod secundum modulum $5$ quadrato $1$ congruum, invenitur quadratum $9$, cui secundum $25$ est congruum, $16$, cui secundum $125$ congruum etc.

\setcounter{satz}{102}
\addtocounter{satz}{-1}
\textbf{102.}

Quod vero attinet ad numeros per $p$ divisibiles, patet, eorum quadrata per $pp$ fore divisibilia, adeoque omnes numeros per $p$ quidem divisibiles, neque vero per $pp$, ipsius $p^n$ fore non-residua. Generaliter vero, si proponitur numerus $p^k A$, ubi $A$ per $p$ non est divisibilis, hi casus erunt distinguendi.

1) Quando $k \geq n$ vel $k=n$, erit $p^k A \equiv 0 \pmod{p^n}$, i.~e. residuum.

2) Quando $k<n$ atque impar, erit $p^k A$ non-residuum.

Si enim esset $p^k A \equiv p^k ss \equiv \pmod{p^n}$, $ss$ per $p^{n-k}$ divisibilis esset, id quod aliter fieri nequit, quam si fuerit $s^2$ per $p^{n-k}$ divisibilis. Tunc vero $ss$ etiam per $p^{n-k+1}$ divisibilis, adeoque etiam (quia $2x+2$ certo non maior quam $n$) $p^k A$, i.~e. $p^{2x+2}A$; sive $A$ per $p$, contra hyp.

3) Quando $k<n$ atque par. Tum $p^k A$ erit residuum vel non-residuum ipsius $p^n$, prout $A$ est residuum vel non-residuum ipsius $p$. Quando enim $A$ est residuum ipsius $p$, erit etiam residuum ipsius $p^{n-k}$. Posito autem $A \equiv aa \pmod{p^{n-k}}$ erit $Ap^k \equiv aap^k \pmod{p^n}$, $aap^k$ vero est quadratum. Quando autem $A$ est non-residuum ipsius $p$, $p^k A$ residuum ipsius $p^n$ esse nequit. Ponatur enim $p^k A \equiv aa \pmod{p^n}$, eritque necessario $aa$ per $p^k$ divisibilis. Quotiens erit quadratum, cui $A$ secundum modulum $p^{n-k}$ adeoque etiam secundum modulum $p$ congruus, i.~e. $A$ erit residuum ipsius $p$ contra hyp.

\setcounter{satz}{103}
\addtocounter{satz}{-1}
\textbf{103.}

Quoniam casum $p=2$ exclusimus, de hoc adhuc quaedam dicenda. Quando numerus $2$ est modulus, numerus quicunque erit residuum, non-residua nulla erunt. Quando vero $4$ est modulus, omnes numeri impares formae $4h+1$ erunt residua, omnes vero formae $4h+3$ non-residua. Tandem quando $8$ aut altior potestas numeri $2$ est modulus, omnes numeri impares formae $8h+1$ erunt residua, reliqui vero, seu ii qui sunt formarum $8h+3$, $8h+5$, $8h+7$ erunt non-residua.

Pars posterior huius propositionis inde clara, quod quadratum cuiusvis numeri imparis, sive sit formae $4k+1$, sive formae $4k-1$, fit formae $8h+1$. Priorem ita probamus.

1) Si duorum numerorum vel summa vel differentia per $2^{n-1}$ est divisibilis, numerorum quadrata erunt congrua secundum modulum $2^n$. Si enim alter ponitur $=a$, erit alter formae $2^{n-1}h + a$, cuius quadratum invenitur $\equiv aa \pmod{2^n}$.

2) Quivis numerus impar, qui ipsius $2^n$ est residuum quadraticum, congruus erit quadrato alicui, cuius radix est numerus impar et $<2^{n-2}$. Sit enim quadratum quodcunque, cui numerus ille congruus, $aa$ atque numerus $a \equiv \pm\tilde{a} \pmod{2^{n-2}}$, ita ut $\tilde{a}$ moduli semissem non superet (art.~4), eritque $aa \equiv \tilde{a}\tilde{a}$. Quare etiam numerus propositus erit $\tilde{a}\tilde{a}$. Manifesto vero tum $a$ tum $\tilde{a}$ erunt impares atque $\tilde{a}<2^{n-2}$.

3) Omnium numerorum imparium ipso $2^{n-2}$ minorum quadrata secundum $2^n$ incongrua erunt. Sint enim duo tales numeri $r$ et $s$, quorum quadrata si secundum $2^n$ essent congrua, foret $(r-s)(r+s)$ per $2^n$ divisibilis (posito $r>s$). Facile vero perspicitur numeros $r-s$, $r+s$ simul per $4$ divisibiles esse non posse, quare si alter tantummodo per $2$ est divisibilis, alter, ut productum per $2^n$ divisibilis fieret, per $2^{n-1}$ divisibilis esse deberet, Q.~E.~A. quoniam uterque $<2^{n-1}$.

4) Quodsi denique haec quadrata ad residua sua minima positiva reducuntur, habebuntur $2^{n-3}$ residua quadratica diversa modulo minora\footnote{Puta quoniam multitudo numerorum imparium infra $2^{n-2}$ est $2^{n-3}$.}, quorum quodvis erit formae $8h+1$. Sed quum praecise $2^{n-3}$ numeri formae $8h+1$ modulo minores exstent, necessario hi omnes inter illa residua reperientur. Q.~E.~D.

Ut quadratum numero dato formae $8h+1$ secundum modulum $2^n$ congruum inveniatur, methodus similis adhiberi potest ut in art.~101: vid. etiam art.~88. --- Denique de numeris paribus eadem valent, quae art.~102 generaliter exposuimus.

\setcounter{satz}{104}
\addtocounter{satz}{-1}
\textbf{104.}

Circa multitudinem valorum diversorum (i.e. secundum modulum incongruorum), quos expressio talis $V = \sqrt{A} \pmod{p^n}$ admittit, siquidem $A$ est residuum ipsius $p^n$, facile e praec. colliguntur haec. (Numerum $p$ supponimus esse primum, ut ante, et brevitatis caussa casum $n=1$ statim includimus).

I. $V$ unum valorem habet pro $p=2$, $n=1$ puta $V=1$; vel si $A$ per $p$ non est divisibilis, duos, quando $p$ est impar, nec non pro $p=2$, $n=2$, puta ponendo unum $=v$, alter erit $=-v$; quatuor pro $p=2$, $n>2$, scilicet ponendo unum $=v$, reliqui erunt $-v$, $2^{n-1}-v$, $2^{n-1}+v$.

II. Si $A$ per $p$ divisibilis est, neque vero per $p^2$, sit potestas altissima ipsius $p$ ipsum $A$ metiens $p^{2\lambda}$ (manifeste enim ipsius exponens par esse debebit) atque $A=ap^{2\lambda}$. Tunc patet, omnes valores ipsius $V$ per $p^\lambda$ divisibiles esse, et quotientes e divisione ortos fieri valores expr. $V' = \sqrt{a} \pmod{p^{n-2\lambda}}$; hinc omnes valores diversi ipsius $V$ prodibunt, multiplicando omnes valores expr. $V'$ inter $0$ et $p^{n-2\lambda}$ sitos per $p^\lambda$; quare illi exhibebuntur per
\[vp^\lambda,\quad vp^\lambda + p^{n-\lambda},\quad vp^\lambda + 2p^{n-\lambda},\quad \ldots\quad vp^\lambda + (p^\lambda-1)p^{n-\lambda},\]
si $V'$ indefinite omnes valores diversos expr. $V'$ exprimit, ita ut illorum multitudo fiat $p^\lambda$, $2p^\lambda$ vel $4p^\lambda$, prout multitudo horum (per casum I) est $1$, $2$ vel $4$.

III. Si $A$ per $p^{2n}$ divisibilis est, facile perspicietur, statuendo $n=2m$ vel $=2m-1$ prout par est vel impar, omnes numeros per $p^m$ divisibiles, neque ullos alios, esse valores ipsius $V$; quare omnes valores diversi hi erunt $0$, $p^m$, $2p^m$, \ldots, $(p^{n-m}-1)p^m$, quorum multitudo $p^{n-m}$.

\subsection*{Criterium generale, utrum numerus datus numeri primi dati residuum sit an non-residuum.}

\setcounter{satz}{105}
\addtocounter{satz}{-1}
\textbf{105.}

Superest casus, ubi modulus $m$ e pluribus numeris primis compositus est. Sit $m = abc\ldots$, designantibus $a$, $b$, $c$ etc. numeros primos diversos aut primorum diversorum potestates, patetque statim, si $n$ sit residuum ipsius $m$, fore etiam $n$ residuum singulorum $a$, $b$, $c$ etc., adeoque $n$ certo non-residuum ipsius $m$ esse, si fuerit non-residuum ullius e numeris $a$, $b$, $c$ etc. Vice versa autem, si $n$ singulorum $a$, $b$, $c$ etc. residuum est, etiam residuum producti $m$ erit. Supponendo enim, $n \equiv A^2$, $B^2$, $C^2$ etc. sec. mod. $a$, $b$, $c$ etc. resp., patet, si numerus $N$ ipsis $A$, $B$, $C$ etc. sec. mod. $a$, $b$, $c$ etc. resp. congruus eruatur (art.~32), fore $n \equiv NN$ secundum omnes hos modulos adeoque etiam secundum productum $m$. --- Quum facile perspiciatur, hoc modo e combinatione cuiusvis valoris ipsius $A$ sive expr.~$\sqrt{n} \pmod{a}$ cum quovis valore ipsius $B$ cum quovis valore ipsius $C$ etc. oriri valorem ipsius $N$ sive expr.~$\sqrt{n} \pmod{m}$, nec non e combinationibus diversis produci diversos $N$, et e cunctis cunctos: multitudo omnium valorum diversorum ipsius $N$ aequalis erit producto e multitudinibus valorum ipsorum $A$, $B$, $C$ etc., quas determinare in art. praec. docuimus. Porro manifestum est, si unus valor expressionis $\sqrt{n} \pmod{m}$ sive ipsius $N$ fuerit notus, hunc simul fore valorem omnium $A$, $B$, $C$ etc.; et quum hinc per art. praec. omnes reliqui valores harum quantitatum deduci possint, facile sequitur, ex uno valore ipsius $N$ omnes reliquos obtineri posse.

\textsc{Ex.} Sit modulus $315$, cuius residuum an non-residuum sit $46$, quaeritur. Divisores primi numeri $315$ sunt $3$, $5$, $7$, atque numerus $46$ residuum cuiusvis eorum quare etiam ipsius $315$ erit residuum. Porro, quia $46 \equiv 1$, et $= 64 \pmod{9}$; $\equiv 1$ et $\equiv 16 \pmod{5}$; $\equiv 4$ et $\equiv 25 \pmod{7}$, inveniuntur radices quadratorum, quibus $46$ secundum modulum $315$ congruus, $19$, $26$, $44$, $89$, $226$, $271$, $289$, $296$.

\setcounter{satz}{106}
\addtocounter{satz}{-1}
\textbf{106.}

Ex praecedentibus colligitur, si tantummodo semper dignosci posset, utrum numerus primus datus numeri primi dati residuum sit an non-residuum, omnes reliquos casus ad hunc reduci posse. Pro illo itaque casu criteria certa omni studio nobis erunt indaganda. Antequam autem hanc perquisitionem aggrediamur, criterium quoddam exhibemus ex Sect. praec. petitum, quod quamvis in praxi nullum fere usum habeat, tamen propter simplicitatem atque generalitatem memoratu dignum est.

Numerus quicunque $A$ per numerum primum $2m+1$ non divisibilis, huius primi residuum est vel non-residuum prout $A^m \equiv +1$ vel $\equiv -1 \pmod{2m+1}$.

Sit enim pro modulo $2m+1$, $A$ index $=a$, in systemate quocunque; eritque $a$ par, quando $A$ est residuum ipsius $2m+1$, impar vero quando $A$ non-residuum. At numeri $A^m$ index erit $ma$, i.e. vel $\equiv 0 \pmod{2m}$, vel $\equiv m \pmod{2m}$, prout $a$ par vel impar. Hinc denique $A^m$ in priori casu erit $+1$, in posteriori vero $-1 \pmod{2m+1}$. V. artt.~57,~62.

\textsc{Ex.} $3$ ipsius $13$ est residuum, quia $3^6 \equiv +1 \pmod{13}$; $2$ vero ipsius non-residuum, quoniam $2^6 \equiv -1 \pmod{13}$.

At quoties numeri examinandi mediocriter sunt magni, hoc criterium ob calculi immensitatem prorsus inutile erit.

\subsection*{Disquisitiones de numeris primis quorum residua aut non-residua sint numeri dati.}

\setcounter{satz}{107}
\addtocounter{satz}{-1}
\textbf{107.}

Facillimum quidem est, proposito modulo, omnes assignare numeros, qui ipsius residua sunt vel non-residua. Scilicet si ille numerus ponitur $=m$, determinari debent quadrata, quorum radices semissem ipsius $m$ non superant, sive etiam numeri bis quadratis secundum $m$ congrui (ad praxin methodi adhuc expeditiores dantur), tuncque omnes numeri horum alicui secundum $m$ congrui, erunt residua ipsius $m$, omnes autem numeri nulli istorum congrui erunt non-residua. At quaestio inversa, proposito numero aliquo, assignare omnes numeros, quorum ille sit residuum vel non-residuum, multo altioris est indaginis. Hoc itaque problema, a cuius solutione illud quod in art. praec. nobis proposuimus pendet, in sequentibus perscrutabimur, a casibus simplicissimis inchoantes.

\subsection*{Residuum $-1$.}

\setcounter{satz}{108}
\addtocounter{satz}{-1}
\textbf{108.}

\textsc{Theorema.} Omnium numerorum primorum formae $4n+1$, $-1$ est residuum quadraticum, omnium vero numerorum primorum formae $4n+3$ non-residuum.

\textsc{Ex.} $-1$ est residuum numerorum $5$, $13$, $17$, $29$, $37$, $41$, $53$, $61$, $73$, $89$, $97$ etc., e quadratis numerorum $2$, $5$, $4$, $12$, $6$, $9$, $23$, $11$, $27$, $34$, $22$ etc. respective oriundum; contra non-residuum est numerorum $3$, $7$, $11$, $19$, $23$, $31$, $43$, $47$, $59$, $67$, $71$, $79$, $83$ etc.

Mentionem huius theor. iam in art.~64 fecimus. Demonstratio vero facile ex art.~16 petitur. Etenim pro numero primo formae $4n+1$ est $(-1)^{2n} \equiv +1$, pro numero autem formae $4n+3$ habetur $(-1)^{2n+1} \equiv -1$. Convenit haec demonstratio cum ea, quam l.~c. tradidimus. Sed propter theorematis elegantiam atque utilitatem non superfluum erit, alio adhuc modo idem ostendisse.

\setcounter{satz}{109}
\addtocounter{satz}{-1}
\textbf{109.}

Designemus complexum omnium residuorum numeri primi $p$, quae ipso $p$ sunt minora, excluso residuo $0$, per literam $C$, et quoniam horum residuorum multitudo semper $=\frac{p-1}{2}$, manifestum est, eam fore parem, quoties $p$ sit formae $4n+1$, imparem vero, quoties $p$ sit formae $4n+3$. Dicantur, ad instar art.~77, ubi de numeris in genere agebatur, residua \emph{socia} talia, quorum productum $\equiv 1 \pmod{p}$; manifeste enim si $r$ est residuum, etiam $\frac{1}{r} \pmod{p}$ residuum erit. Et quoniam idem residuum plura socia inter residua $C$ habere nequit, patet, omnia residua $C$ in classes distribui posse, quarum quaevis bina residua socia contineat. Iam perspicuum est, si nullum residuum daretur, quod sibi ipsi esset socium, i.~e. si quaevis classis bina residua inaequalia contineret, omnium residuorum numerum fore duplum numeri omnium classium; quodsi vero aliqua dantur residua sibi ipsis socia i.~e. aliquae classes, quae unicum tantum residuum aut, si quis malit, idem residuum bis continent, posita harum classium multitudine $=a$, reliquarumque multitudine $=b$; erit omnium residuorum $C$ numerus $=a+2b$. Quare quando $p$ est formae $4n+1$, erit $a$ numerus par; quando autem $p$ est formae $4n+3$, erit $a$ impar. At numeri ipso $p$ minores alii, quam $1$ et $p-1$, sibi ipsis socii esse nequeunt (vid. art.~77); priorque $1$ certo inter residua occurrit; unde in priori casu $p-1$ (seu quod hic idem valet, $-1$) debet esse residuum, in posteriori vero non-residuum; alias enim in illo casu foret $a=1$, in hoc autem $=2$, quod fieri nequit.

\setcounter{satz}{110}
\addtocounter{satz}{-1}
\textbf{110.}

Etiam haec demonstratio ill. Euleri debetur, qui et priorem primus invenit V. Opusc. Anal. T.~I.~p.~125. --- Facile quisquis videbit eam similibus principiis innixam esse, ut demonstratio nostra secunda theor. Wilsoniani art.~77. Si vero hoc theorema supponere velimus, facilius adhuc demonstratio exhiberi poterit. Scilicet inter numeros $1$, $2$, $3$, \ldots, $p-1$ erunt $\frac{p-1}{2}$ residua quadratica ipsius $p$ totidemque non-residua; quare non-residuorum multitudo erit par, quando $p$ est formae $4n+1$; impar, quando $p$ est formae $4n+3$. Hinc productum ex omnibus numeris $1$, $2$, $3$, \ldots, $p-1$ in priori casu erit residuum, in posteriori non-residuum (art.~99). At productum hoc semper $\equiv -1 \pmod{p}$; adeoque etiam $-1$ in priori casu residuum, in posteriori non-residuum erit.

\setcounter{satz}{111}
\addtocounter{satz}{-1}
\textbf{111.}

Si itaque $r$ est residuum numeri alicuius primi formae $4n+1$, etiam $-r$ huius primi residuum erit; omnia autem talis numeri non-residua, etiam signo contrario sumta non-residua manebunt\footnote{Quando igitur de numero quocunque loquemur quatenus numeri formae $4n+1$ residuum vel non-residuum est, ipsius signum omnino negligere sive etiam signum anceps $+$ ipsi tribuere poterimus.}. Contrarium evenit pro numeris primis formae $4n+3$, quorum residua quando signum mutatur, non-residua fiunt et vice versa, vid. art.~98.

Ceterum facile ex praecedentibus derivatur regula generalis: $-1$ est residuum omnium numerorum, qui neque per $4$ neque per ullum numerum primum formae $4n+3$ dividi possunt; omnium reliquorum non-residuum. V. artt.~103 et~105.

\subsection*{Residua $+2$ et $-2$.}

\setcounter{satz}{112}
\addtocounter{satz}{-1}
\textbf{112.}

Progredimur ad residua $+2$ et $-2$.

Si ex tabula II colligimus omnes numeros primos, quorum residuum est $+2$, hos habebimus: $7$, $17$, $23$, $31$, $41$, $47$, $71$, $73$, $79$, $89$, $97$. Facile autem animadvertitur, inter hos numeros nullos inveniri formarum $8n+3$ et $8n+5$.

Videamus itaque, num haec inductio ad certitudinem evehi possit.

Primum observamus, quemvis numerum compositum formae $8n+3$ vel $8n+5$ necessario factorem primum alterutrius formae $8n+3$ vel $8n+5$, involvere; manifesto enim e solis numeris primis formarum $8n+1$, $8n+7$, alii numeri quam qui sunt formae $4n+1$ vel $4n+3$, componi nequeunt. Quodsi itaque inductio nostra generaliter est vera, nullus omnino numerus formae

\end{document}
